% Latin 11195, batch 1 — Gallica views 1–20.
% Pass: Opus 5 (claude-opus-5), 2026-09-19 — first pass, unchecked against
% the views by a human.
% The pass was resumed after an interruption: views 3 to 9 L had been drafted
% before the break and are taken over as they stood, except for two notes on
% views 4 R and 7 R that put the running marginal figures in the outer margin
% when on a right page they stand in the inner one. Views 1, 2 and 9 R to 20
% were read from the mirror in this pass and drafted view by view before the
% file was assembled.
%
% What the views are. Greyscale scans, landscape, 5727 × 3897 (5881 × 4051 for
% view 1), each showing the volume lying open: a left page and a right page
% with the gutter near the middle, its position shifting from view to view
% (about x = 2668 at view 4, x = 2920 at view 11). Each view was split at the
% gutter and read as two pages; in the file each \page{N} holds the left page
% then the right, and a \note{} says where the right page begins. \page{N} is
% always the Gallica view number.
%
% Notation. The period's sign of equality is a double vertical stroke, « ∥ »,
% and is set \parallel in mathematics throughout; it is never modernised to
% «=». The same stroke serves for a proportion in the marginal calculation of
% view 9 R. The species notation is Viète's as the Paris school wrote it:
% A the unknown, B C D F the given sides, x the coefficient of the next-to-
% highest degree, and the homogeneous constants named by their dimension —
% « z pl. » (zetetic plane), « T Sol. » (solid), « S pl.pl. » (plano-plane),
% « P Sol. », « G pl. », « B pl. ». Abbreviations are kept as written with
% their marks — ōes, ōib⁹ (= -bus), ē (= est), ꝑ, æquatiōe, Atq;, & — and are
% expanded only in a \note{}. The long s is set s.
%
% Blank views, absent from the file. View 1: the inside of the upper board,
% with the printed shelf label « LATIN 11195 », and the recto of the flyleaf,
% blank but for the old shelfmark « Suppl.t lat. 218. A. » at the top outer
% corner. The right pages of views 14 and 15 are blank too, and are said to be
% so in the notes of those views; the left page of view 2 (verso of the
% flyleaf) likewise. The gap in the numbering is the record.
%
% What the twenty views hold. One text, in one regular seventeenth-century
% cursive, in Latin: a treatise « De Recognitione Æquationum » in the Viète
% school's manner, on how an equation is put together out of its roots and how
% it is to be read back — chapter by chapter, from quadratics up.
%
%   v. 2 R     The title leaf: « De Recognitione Æquationum », three words in
%              the middle of the upper third, with an ink « 171 » above, cut
%              by the binder's knife and in another hand.
%   v. 3 R     The treatise opens, paginated « 1 » by the writer. The general
%              doctrine: what recognition is, why an equation is explicable of
%              as many sides as went into it, and the three kinds of side —
%              true, « positiva supra » and « positiva infra » (which stand
%              for the negative roots), and fictitious ones, whose being the
%              sole explanation of an equation marks the question impossible.
%   v. 4 R     « De constitutione Æquationum Quadraticarum », with Caput 1
%              (both sides above), Caput 2 (both below) written side by side
%              in two columns, and Caput 3 (one above, one below).
%   v. 5       The three propositions of Caput 3; Caput 4, on the quadratic
%              that subsists by itself, « z pl. + A² ∥ 0 »; then the reading
%              of the canonical forms and the beginning of the doctrine of
%              « determinatio ».
%   v. 6–8 L   « De Constitutione Æquationum Cubicarum »: six capita, with the
%              products written out and each constitution interpreted into its
%              canonical form.
%   v. 8 R–11  « De Determinationibus Æquationum Cubicarum »: the case of two
%              equal roots reduced by successive comparisons; the greatest
%              parallelepiped applied to a given plane and deficient by a
%              cubic figure; and, twice named, lemmas « demonstratum a D. Rob. »
%              — Roberval — with a marginal note in French sending the reader
%              to « son traité de mechanique des plans inclinez » (v. 9 R). The
%              determinations end at view 11 R, the lower half of which is
%              blank.
%   v. 12–17   « De constitutione Æquationum Quadratoquadraticarum », ten capita
%              announced; Caput 1 (four sides above), Caput 2 (three above, one
%              below, seven propositions), Caput 3 (two and two, thirteen
%              propositions), Caput 4 (three below, one above, seven), Caput 5
%              (all four below).
%   v. 18–20   Caput 6, the quadrato-quadratic born of two plane equations;
%              Caput 7, of a plane and two laterals, with eleven propositions;
%              Caput 8, of a cubic and a lateral, running to its twentieth
%              proposition at the foot of view 20 R. The text continues at
%              view 21, in batch 2.
%
% Two leaves out of the run. The leaf paginated 20–21 carries, pinned to its
% top outer corner, two small slips of paper. In view 13 the first of them
% covers the page number; views 14 and 15 photograph the same blank verso with
% the slips lifted one after the other, and the number « 20 » then reads on
% the projecting edge of the leaf. Each slip is an addition to the text and is
% transcribed under the view that shows it: the first is headed « p. 21. » and
% supplies a sentence on the two affirmed rectangles; the second, headed
% « 19a », asks how one is to tell whether a given cubic belongs to the fifth
% chapter or the third.
%
% A second hand in the margins. A finer pen, paler, has copied the canonical
% form of each constitution into the left margin beside it, and numbered its
% own remarks 1, 3, 4, 5, 6 … up to 29 down the same margin (on a right page
% that margin is the inner one). It also wrote the numerical check of the
% lemma at view 9 R, the word « minorem » against an underlined « maiorem »
% at view 11 R, and small star-like marks under some signs. Where its copied
% form and the line disagree in a sign — view 12 R prop. 2, view 20 L prop. 5
% — both are given as they stand and a \note{} says so.
%
% Numbers on the leaves. The ink figures at the outer top corners are the
% writer's own pagination, 1 (v. 3 R) to 31 (v. 20 R), one per page; they are
% recorded in the \note{} of each view and never in \folio{}. The title leaf
% carries a different ink « 171 » at its head, the flyleaf the old shelfmark
% « Suppl.t lat. 218. A. », and view 3 the old shelfmark again, « Suppl.t Lat.
% no 218. » at the head and « Suppl.t l. 218. + I. » at the foot, with a « 26 »
% on the blank left page and the round stamp « Bibliothèque impériale MSS ».
% At views 19 and 20 a bold « 33 » shows on the projecting top edge of a later
% leaf: it is that leaf's pagination, seen in advance. No pale pencil series,
% one number per leaf at the top right of rectos, was found anywhere in the
% batch, so no \folio{} is written.

\documentclass[11pt,a4paper]{article}
% The preamble shared by every transcription of the mathematicians' manuscripts in Gallica.
%
% It rests on one idea: **uncertainty compounds, it does not resolve.** A
% transcription that smooths over a doubtful word has destroyed the only thing
% it was made for — a reader must be able to tell, at every point, what was on
% the page and what was supplied. The apparatus macros below are therefore the
% critical apparatus, not decoration.
%
% The same commands are recognised by `scripts/render.mjs`, which produces the
% site's reading view, and by `scripts/tei.mjs`. A macro added here must be
% added there too, or rendering fails loudly — which is the wanted behaviour.
%
% Comments here are English, like the rest of the repository. The documents
% this preamble sets are French, and that is deliberate: see the skills.

\ProvidesPackage{germain}[2026/09/15 Transcriptions of mathematicians' manuscripts in Gallica]

\usepackage{amsmath,amssymb,amsthm}
\usepackage{mathrsfs}
% Kept from the parent preamble so the renderer's diagram support stays
% available; these manuscripts draw no commutative diagrams, and nothing here uses it.
\usepackage{tikz-cd}
\usepackage[english,french]{babel}
\usepackage{fontspec}

% --- Greek ------------------------------------------------------------------
% The text face stays fontspec's default, Latin Modern, which has no Greek:
% XeTeX drops every glyph it lacks with a line in the log and nothing on the
% page, and the Greek words the manuscripts quote vanished from the PDFs. So
% the Greek and Greek Extended blocks get a XeTeX character class of their own,
% and entering or leaving a run of it switches the family to CMU Serif — the
% same Computer Modern design, with polytonic Greek — and back. Only the
% family changes: a Greek word in \textit or \textbf keeps its shape and series.
% The transitions to and from every other class carry the tokens that class
% has towards an ordinary letter, so French spacing before « ; » or inside
% « » still holds after a Greek word. No grouping, so no brace or \endgroup
% can come out unbalanced; maths is left alone, since XeTeX inserts nothing
% there. Kept identical in `exercices.sty`.
\makeatletter
\newfontfamily\germainGreekfont{cmun}[
  Extension=.otf, NFSSFamily=germaingreek,
  UprightFont=*rm, ItalicFont=*ti, SlantedFont=*sl,
  BoldFont=*bx, BoldItalicFont=*bi, BoldSlantedFont=*bl]
\newXeTeXintercharclass\germain@greek
\def\germain@greekrange#1#2{%
  \@tempcnta=#1\relax
  \loop
    \XeTeXcharclass\@tempcnta=\germain@greek
    \ifnum\@tempcnta<#2\advance\@tempcnta\@ne\repeat}
\germain@greekrange{"0370}{"03FF}
\germain@greekrange{"1F00}{"1FFF}
\def\germain@greekprev{\rmdefault}
\def\germain@greekin{%
  \edef\germain@greekprev{\f@family}\fontfamily{germaingreek}\selectfont}
\def\germain@greekout{\fontfamily{\germain@greekprev}\selectfont}
\def\germain@greekjoin{%
  \XeTeXinterchartokenstate=\@ne
  \@tempcnta=\z@
  \loop
    \ifnum\@tempcnta=\germain@greek\else
      \XeTeXinterchartoks\germain@greek\@tempcnta=\expandafter{%
        \expandafter\germain@greekout\the\XeTeXinterchartoks\z@\@tempcnta}%
      \XeTeXinterchartoks\@tempcnta\germain@greek=\expandafter{%
        \the\XeTeXinterchartoks\@tempcnta\z@\germain@greekin}%
    \fi
    \ifnum\@tempcnta<4095\advance\@tempcnta\@ne\repeat}
% babel-french sets its own classes, and saves and restores the interchar
% state, each time a language is selected: join again after it does.
\addto\extrasfrench{\germain@greekjoin}
\addto\extrasenglish{\germain@greekjoin}
\AtBeginDocument{\germain@greekjoin}
\makeatother

\usepackage{geometry}
\geometry{a4paper,margin=25mm,marginparwidth=28mm}
\usepackage{marginnote}
\usepackage{xcolor}
\usepackage[normalem]{ulem}
\setlength{\emergencystretch}{3em}
\usepackage{hyperref}
\hypersetup{colorlinks=true,linkcolor=black,urlcolor=black,pdfborder={0 0 0}}

% --- Batch metadata ---------------------------------------------------------
% Taken as they stand by the rendering script, which shows them at the head of
% the reading view. They fix what the file claims to cover. The macro names are
% the parent project's, so that the scripts stay shared; \folder here means the
% volume's slug — `fr-9115` — and \pages counts Gallica views.

\newcommand{\folder}[1]{\gdef\grothFolder{#1}}
\newcommand{\batch}[1]{\gdef\grothBatch{#1}}
\newcommand{\pages}[2]{\gdef\grothFirst{#1}\gdef\grothLast{#2}}
\newcommand{\foldertitle}[1]{\gdef\grothFoldertitle{#1}}

% The holder's own shelfmark, printed as they write it: « Français 9115 ».
\newcommand{\shelfmark}[1]{\gdef\grothShelfmark{#1}}

% The dating, copied verbatim from the holder's catalogue — never inferred
% here. « XIXe siècle » is the BnF's; a pass that dates a leaf from its content
% says so in a \note{}, as its own inference.
\newcommand{\dating}[1]{\gdef\grothDating{#1}}

% The status notice, printed in the title block and repeated as a diagonal
% watermark on every page of the PDF. The manuscripts are in the public
% domain; what this line declares is not a right but a quality — an
% unverified first pass by a machine. The skills write the line; a file
% without it gets no watermark, so leaving it out is visible in the output.
\newcommand{\watermark}[1]{\gdef\grothWatermark{#1}}

\gdef\grothFolder{?}\gdef\grothBatch{}\gdef\grothFirst{?}\gdef\grothLast{?}%
\gdef\grothFoldertitle{}\gdef\grothShelfmark{}\gdef\grothDating{}\gdef\grothWatermark{}

% --- The résumé -------------------------------------------------------------
\newenvironment{resume}
  {\small\setlength{\parskip}{0.45em}}
  {\par\medskip\hrule\medskip}

% --- Keywords ---------------------------------------------------------------
% In English, closing the modernised reading's résumé: the modern vocabulary
% under which to search for what the leaves construct. The manifest extracts
% them to tag the volume — this line is the single source of the tags.
\newcommand{\keywords}[1]{\par\smallskip\noindent
  {\footnotesize\textsc{Keywords} --- \textit{#1}}\par}

% --- The critical apparatus -------------------------------------------------

% The Gallica view number, in the margin. This is the anchor that lets a
% disagreement over a formula be settled by looking at *one* image rather than
% a volume of 750. The site uses it to turn the facsimile as the transcription
% is scrolled. A view is one scanned image: usually one side of a leaf.
\newcommand{\page}[1]{%
  \marginnote{\footnotesize\color{gray!70}\textsf{v.\,#1}}%
  \label{page:#1}\ignorespaces}

% The BnF's pencilled foliation, where the leaf carries it and a pass can read
% it: « 198r ». This is what the literature cites — « ff. 198r–208v » — and
% the one line that turns a citation into a view. Recorded, never inferred:
% a folio a pass cannot read is not written.
\newcommand{\folio}[1]{%
  \marginnote{\footnotesize\color{gray!50}\textsf{f.\,#1}}\ignorespaces}

% The modernised reading's coarser counterpart to \page{}: which views a
% section is drawn from.
\newcommand{\pagerange}[2]{%
  \marginnote{\footnotesize\color{gray!70}\textsf{v.\,#1--#2}}%
  \ignorespaces}

% Illegible: never guessed.
\newcommand{\ill}{\textcolor{red!60!black}{[\,\ldots\,]}}

% A doubtful reading: the word is given, and flagged as uncertain.
\newcommand{\uncertain}[1]{\uline{#1}}

% An editorial addition — a missing letter, an implied word.
\newcommand{\add}[1]{\textcolor{blue!50!black}{[#1]}}

% Struck out by the writer. What was crossed out is part of the document.
\newcommand{\struck}[1]{\sout{#1}}

% The transcriber's note, on the state of the page or the mathematics.
\newcommand{\note}[1]{\par\smallskip{\small\color{gray!60!black}\textit{[#1]}}\par\smallskip}

% A marginal note of hers, distinct from the body of the text.
\newcommand{\marginal}[1]{\par\smallskip\noindent\hspace{1em}%
  {\small\color{gray!60!black}#1}\par\smallskip}

% --- Title ------------------------------------------------------------------
% The title block is French, like the documents themselves.

\AddToHook{shipout/background}{%
  \ifx\grothWatermark\empty\else
    \begin{tikzpicture}[remember picture, overlay]
      \node[rotate=52, scale=3.4, align=center, text=gray!18,
            font=\sffamily\bfseries]
        at (current page.center) {\grothWatermark};
    \end{tikzpicture}%
  \fi}

\AtBeginDocument{%
  \begin{center}
    {\large\bfseries Manuscrits de mathématiciens (Gallica) ---
      \ifx\grothShelfmark\empty\grothFolder\else\grothShelfmark\fi}\\[2pt]
    {\itshape\grothFoldertitle}\\[4pt]
    {\small vues \grothFirst--\grothLast
      \ifx\grothBatch\empty\else\ (lot \grothBatch)\fi}\\[2pt]
    \ifx\grothDating\empty\else
      {\small\color{gray!60!black}Datation du catalogue : \grothDating}\\[2pt]
    \fi
    \ifx\grothWatermark\empty\else
      {\small\bfseries\color{red!45!gray}\grothWatermark}\\[2pt]
    \fi
    {\footnotesize\color{gray!60!black}Transcription établie d'après les images IIIF de Gallica.
     Source gallica.bnf.fr / Bibliothèque nationale de France.\\
     Les lectures incertaines sont soulignées, les illisibles notées [\,\ldots\,],
     les ajouts entre crochets ; \textsf{v.} numérote les vues de Gallica,
     \textsf{f.} la foliotation au crayon de la BnF.}
  \end{center}
  \vspace{1em}}

\endinput


\folder{latin-11195}
\shelfmark{Latin 11195}
\batch{1}
\pages{1}{20}
\dating{}
\watermark{Lecture automatique\\non vérifiée}
\foldertitle{Mélanges de physique et de mathématiques, comprenant des opuscules et des lettres de P. de Roberval, Huggens, Torricelli, Firmat et Fr. Herman Flayder. Latin 11195}

\begin{document}

\page{2}
\note{Double page. La page de gauche, verso de la garde, est blanche. La page de droite est le feuillet de titre : trois mots au tiers supérieur, de la main du traité, et au-dessus, à l'encre et rogné par le couteau du relieur, le nombre « 171 », d'une autre main que la pagination du traité.}

\section*{De Recognitione Æquationum}

\page{3}
\note{Double page. La page de gauche, dos du feuillet de titre, est blanche ; elle porte au milieu, à l'encre, le nombre « 26 ». Page de droite : début du traité, paginée « 1 » à l'encre en haut à l'angle extérieur. En haut, d'une main de bibliothécaire, l'ancienne cote « Suppl.\textsuperscript{t} Lat. n\textsuperscript{o} 218. » ; au pied de la page, de la même main et soulignée, « Suppl.\textsuperscript{t} l. 218. + I. ». Dans la marge de gauche, le cachet rond « Bibliothèque impériale MSS ».}

\section*{De Recognitione Æquationum}

Ideo is tractatus d\textsuperscript{r} de æquationum recognitiōe quia in ipso investigantur modj quibus æquationum variæ constitutiones dignoscantur, seu per quos explicentur. Cum enim ob multiplices tum ductiones, tum aplicatiōes, tū alias analyseos operatiōes dictæ æquationes sub diversis formulis delitescant, ita ut earū radices, seu latera, his involutæ ambagibus difficulter se prodant, huius partis ē eas radices seu latera exhibere.

Præterea cū ōis æquatio vel ꝑ se subsistat, vel ex aliā componatur, hæc recognitio erit quædam veluti resolutio, eorū quæ multiplicatiōis ope producta fuerant, \& sic ōia quib\textsuperscript{9} potestas in altiorē gradū in quo reperitur ascenderit, facile dignoscentur, æquatio enim in se habet eas ōes radices ex quibus fuit composita, ac proinde de iisdem ōib\textsuperscript{9} est explicabilis. Atq; inde sequitur eo magis implicatam æquationem evenire, ac de plurib\textsuperscript{9} laterib\textsuperscript{9} ē explicabilē, quo gradus potestatis fuerit altior; nam v.g. æquatio quadratica cū vel à se subsistat, vel ex ductu duarū lateralium tantū possit procreari; de iisdē etiam erit solummodo explicabilis, sed cubica cū \& à se subsistat, \& à tribus lateralibus, \& à quadratica \& laterali componatur, de iisdem ōib\textsuperscript{9} explicabitur. Et sic in infinitū ubi potestas in altiori gradu reperiretur. Qua in re veteres analystæ densis offusi tenebris plane cæcutiebant, vix enim unicū poterant latus exhibere, atq; inde etiam patet huius methodj elegantia, quæ in eo, ut diximus, est præcipua ut ōia omnino latera ex quibus quævis æquatio composita fuerit exhibeat.

Horū autem laterū quædam habent veram significatiōem, quædam intelliguntur ut nihilo minora, quædam denique nihil veri in se continent sed sunt fictitia, quæ quidem fictitia si sola æquatiōis explicatiōi inserviant, propositæ quæstionis impossibilitatem denotabunt.

Vera latera dividuntur in positiva supra, seu in ea quæ diximus veram habere significatiōem, \& positiva infra, seu quæ tanquam nihilo minora intelliguntur. Positiva suprà sunt quæ aliquid exprimunt præter vel supra nihilum. v.g. si $+A$, sit positivum supra, eiusq; valor sit 4, intelligetur valorem ipsius $A$, esse quatuor præter vel supra nihilū. Positiva vero infrà quamvis concipiantur ut minora nihilo, \& is conceptus videatur maxime fictitius, natura enim nihil producit infrà nihilū, nō sunt tamen maxime fictitia, mutatis enim signis emergunt positiva suprà, quemadmodū \& positiva suprà etiam \& mutatis signis, ad naturam positivorū infra reducuntur, nam idem $A$, quem diximus cū signo $+$ esse valoris 4 supra nihilū, si sit cū signo $-$ erit valoris 4 infra nihilum.

\page{4}
\note{Double page ; la page de gauche est paginée « 2 » à l'encre en haut à l'angle extérieur, celle de droite « 3 ». Dans la marge extérieure de la page de gauche, une série de petits chiffres de la même encre, « 1. », « 3. », « 4. », « 5. », posés en regard des lignes ; ils sont rognés en partie par le couteau du relieur.}

pariq; ratiōe si positivum infrà fuerit cū signo $-$ fit positivum suprà, cū enim negatio negat ipsum \& aliquid infra nihilū, asserit esse aliquid supra nihilum. Sæpius autem occurrit has signorū mutatiōes miræ esse utilitatis, quæ facilius percipientur in tractatu de locis.

Hæc vero latera positiva infrà, ideo excogitata sunt, quia cū in æquationū propositiōe cuiusdam lineæ quantitas determinanda est, quandoq; expedit (\& hac ratiōe variæ \& implicitæ vitantur anfractus) ut hæc linea cuius quantitas investigatur, alteri cognitæ comparetur, \& cuius comparatiōe hæc quantitas ignota elicienda est, aliqualem eius partem solummodo assumere, quæ quidem pars forsan erit minor nihilo, cū ipsa tamen tota incognita possit esse minor ea cognitā cui comparatur, ut sic manifestū est id quod ab ea incognitā rescindendū esse minus nihilo, nihilominus hâc divisiōe ad ipsi cognitiōem devenitur, quā tandem cognitā quod ablatū est ei restituitur, cū tamen hæc pars minor nihilo suas vices egerit in compositiōe huiusmodj æquatiōis, consequens est ut in ea dignoscatur, ac proinde eadem æquatio de ea explicetur.

Duo supersunt notanda in his lateribus quæ sunt positiva infra. 1\textsuperscript{o} latus quod eiusmodj naturam sortitus est, efficere ut ōia quæ sub ipso affecta sunt, sint etiam infra.

Secundò hæc latera in hac methodo, sicut \& positiva supra numquam mutare signa, quæ ipsis præfixa sunt, sive plus, sive minus, sed cū duo latera diversæ naturæ ducuntur alterū in alterū, v.g. cū unū supra cui præficitur signū $+$, ducitur in alterū infrà, cui signū $-$, productū quod oritur notari quidem signo $-$, \& fieri positivum infra, sed per æquivalentiam (ut sic loquar) ita esse accipiendū ac si esset suprà, ut patet ex iis quæ dicta sunt, sic \& latus suprà cui præfixū est signū $-$, si ducatur in latus infrà cui præfixū sit signū $+$, productū quod inde oritur notari quidem signo $-$, \& esse suprà, sed per æquivalentiam ita esse accipiendū ac si esset infrà, sit v.g. quædam æquatio in cuius compositiōe sit hæc lateralis $B + A \parallel 0$, quæ inde oritur æquatio erit explicabilis de uno latere positivo infrà scilicet $+A$ infra, quod per æquivalentiam idem est ac $-A$ supra.

His præmissis æquatio (quæ est aggregatū pluriū terminorū partim cognitorū partim vero ignotorū) in hac methodo sic est constituenda, ut quantitas ignota sit ex eadem parte cū cognitā, \& utraq; pars æquatiōis sic in unū redacta æquetur nihilo, quod quam sit utile patebit in sequentib\textsuperscript{9}. Et tunc affirmata æqualia erunt negatis, sic autem instituta æquatiōe, latera seu radices sunt inquirendæ, sive sint earū quæ sunt suprà, vel infrà, vel etiam fictitiæ, ut \& de omnibus iis explicetur, sit in exemplū sequens æquatio.
\[ \begin{array}{l} 4 - 1\,\mathrm{R} \parallel 0 \\ 3 - 1\,\mathrm{R} \parallel 0 \\ 5 - 1\,\mathrm{R} \parallel 0 \end{array} \qquad \text{productum est } 60 - 47\,\mathrm{R} + 12\,q - 1\,C \parallel 0 \]

\note{Page de droite, paginée « 3 ». Un chiffre « 6. » dans la marge de gauche, qui est ici la marge intérieure, en regard du Caput primum.}

\uncertain{\&} vero quidem æquatio de tot lateribus explicabitur, quot in ipsius compositiōem admissa sunt, sicut in hoc exemplo de tribus lateribus, 4. 3. et 5. quæ quidem latera ut videre est, sunt positiva suprà cū enim $4 - \mathrm{R} \parallel 0$ erit per antithesim $4 \parallel 1\,\mathrm{R}$, secus autem, si hanc habuissemus æquationem $4 + 1\,\mathrm{R} \parallel 0$. eo enim casu 4 fuisset latus positivum infrà et cæt. quæ ex dictis sunt perspicua.

\subsection*{De constitutione Æquationum Quadraticarum.}

Prælibatis generalibus, quæ ad æquatiōes, et earū radices spectant, sequitur ut de unaquaq; æquatiōis constitutiōe sigillatim inquiramus, ac 1\textsuperscript{o} ordiemur à quadraticâ, omniū enim est minimè composita, nec sub se habet aliam magis simplicem exceptâ laterali, de quâ cū nulla occurrat difficultas nihil etiam dicemus.

Æquatio quadrata est cū potestas quadratū est, quod quidem quadratū vel ex duob\textsuperscript{9} lateribus componitur, vel ex se subsistit, ac proinde et hæc æquatio de duob\textsuperscript{9} lateribus erit explicabilis.

Hæc autem duo latera vel erunt ambo suprà, vel ambo infrà, vel unū suprà, alterū vero infrà. Et hæc inductio est legitima cū numerus variorū est finitus, \& latius in geometricas usitatos solita. Et secundū hos varios casus tot erunt capita in quibus etiam tot propositiones, quot erunt differentes laterū, quæ diversis signis notabuntur, magnitudines.

\subsection*{Caput primum}

Quando ambo latera sunt suprà. Propositio unica.
\[ \begin{array}{l} \text{Sit } B - A \parallel 0 \\ \text{Et } C - A \text{ cuiuscumq; valoris} \\ \text{vel.} \\ C - A \parallel 0 \\ \text{\& } B - A \text{ cuiuscumq; valoris} \\ \text{productū erit} \\ BC - CA - BA + A^{2} \parallel 0 \end{array} \]

\subsection*{Caput secundum}

Quando ambo latera sunt infrà. Propositio unica.
\[ \begin{array}{l} \text{Sit } B + A \parallel 0 \\ \text{Et } C + A \text{ cuiuscumq; valoris} \\ \text{vel.} \\ C + A \\ \text{\& } B + A \text{ cuiuscumq; valoris} \\ \text{productū erit} \\ BC + BA + CA + A^{2} \parallel 0 \end{array} \]
\note{Les deux \emph{capita} sont écrits côte à côte, en deux colonnes séparées par un trait vertical. Dans les deux produits, les deux termes du milieu sont réunis par une accolade devant le signe qui les précède : « $BC - \{CA,\ BA\} + A^2$ », « $BC + \{BA,\ CA\} + A^2$ » ; ils sont ici développés.}

\subsection*{Caput tertium}

Quando unū latus est suprà et alterū infrà.

Ad constitutiōem æquationū huius capitis. Notandū est ob laterū varietates diversas etiam evenire signorū $+$ et $-$ notationes, quoniam enim signa sequuntur naturam maioris terminj, si terminus seu species cui præficitur signū $+$ sit maior, sive si latus suprà sit maius latere infrà, quando interpretabitur hæc æquatio, ut jam videbimus, differentia laterum notanda erit cum signo lateris maioris. Sit Igitur.

\page{5}
\note{Double page ; la page de gauche est paginée « 4 », celle de droite « 5 ». Dans la marge extérieure de la page de gauche, les chiffres « 7 » et « 8 » de la série courante, et les formes canoniques de l'équation traitée, recopiées en regard du texte.}

\subsection*{Propositio prima}

Cum latus suprà maius est latere infrà.
\[ \begin{array}{l} \text{Sit } B - A \parallel 0 \\ \text{Et } C + A \parallel \text{ cuiuscumq; valoris} \\ \text{vel.} \\ C + A \parallel 0 \\ \text{\& } B - A \text{ cuiuscumq; valoris} \\ \text{productum erit} \\ BC + BA - CA - A^{2} \parallel 0 \end{array} \]

\subsection*{Propositio secunda}

Cum latus suprà æquale est lateri infrà.
\[ \begin{array}{l} \text{Sit } B - A \parallel 0 \\ \text{Et } C + A \text{ cuiuscumq; valoris et cæt. ut supra} \\ \text{productum erit} \\ BC - A^{2} \parallel 0 \end{array} \]

\subsection*{Propositio Tertia}

Cum latus infrà maius est latere suprà. Erit idem productum scilicet
\[ BC + BA - CA - A^{2} \parallel 0 \]
\note{Dans les deux produits, « $BC$ » est suivi d'une accolade réunissant $BA$ et $CA$, le $+$ et le $-$ étant écrits l'un au-dessus de l'autre devant l'accolade ; ils sont ici développés.}

\subsection*{Caput Quartum. De Æquatione quadraticâ quæ per se subsistit}

Æquatio dicitur per se subsistere, vel a se esse, quæ à qualibet aliâ per quamcumq; ductiōem produci non potest, quoniam vero inter quadraticas unica est eiusmodj ideo huius capitis erit

\subsection*{Propositio unica}

\marginal{$z\,\mathrm{pl.} + A^{2} \parallel 0$}
\[ z\,\mathrm{pl.} + A^{2} \parallel 0 \]
In quâ æquatiōe evidens est unicū esse latus positivū infrà potentia æquale $z\,\mathrm{pl.}$ nam per antithesim erit $z\,\mathrm{pl.} \parallel -A^{2}$.

\subsection*{Ad Caput primum Et secundum}

Fusius in hac parte de æquationib\textsuperscript{9} quadraticis disserimus de iis quæ ad generatim tum æquationū constitutiōem tum interpretatiōem, ne quid supersit quod in sequentib\textsuperscript{9} pariat difficultatem, cū et ea quæ hic dicuntur iis etiam inserviant.

Constituta igitur æquatiōe ut evidenter fieret ipsius origo, ac proinde radices innotescerent, sic eam Interpretamur In primo capite ut scilicet $BC$ sit $z\,\mathrm{pl.}^{um}$, Et $C + B$ sit $X$ seu summa laterum, Cum Igitur similis occurret æquatio.
\marginal{$z\,\mathrm{pl.} - xA + A^{2}$}
\[ Z\,\mathrm{pl.} - xA + A^{2} \parallel 0 \]
erunt duo latera supra quorū summa est $x$, rectangulū sub ipsis est $z\,\mathrm{pl.}$ et $A$ fit, alterutrū ex lateribus.

Et quæcumq; varietates inciderint eodem semper modo, hæc æquatio poterit explicari sit \struck{enim} v.g.
\[ \frac{F^{\mathrm{sd}}}{G} - \frac{KNA}{M} + \frac{HLA}{G} + A^{2} \parallel 0 \]
$-\dfrac{KNA}{M} + \dfrac{HLA}{G}$ Intelligitur tanquam una linea, quæ vocatur $x$ et eodem redit et æquatio cum superiore similiter æquatiōe secundj cap\textsuperscript{is} sic Interpretabimur
\marginal{$z\,\mathrm{pl.} + xA + A^{2}$}
\[ Z\,\mathrm{pl.} + xA + A^{2} \parallel 0 \]
Et erunt ut patet ex ipsius origine, duo latera infra quorū summa est $x$ rectangulū sub ipsis $z\,\mathrm{pl.}$, et $A$, fit quodlibet ex lateribus.

Æquationibus sic Interpretatis videndū est an determinationj sint obnoxiæ, æquatio vero dicitur determinationj obnoxia, cū de duobus vel pluribus lateribus æqualibus

\note{Page de droite, paginée « 5 ». Le tiers inférieur est resté blanc. Les formes canoniques sont répétées dans la marge extérieure.}

est explicabilis, quæ quidem determinatio in universū duplex est maior scilicet et minor, maior cū ōia latera de quibus potest explicari æquatio, sunt æqualia minor verò cū non ōia latera sunt æqualia, sed quædam duntaxat, In quâ autem æquatiōe detegenda est illa constitutio quæ determinationj sive maiorj sive minorj est obnoxia, in prædictis duab\textsuperscript{9} formulis, quarū prima præcipua est inter quadraticas, determinatio est, cū $x$, bifariam secatur et utraq; medietas sit $A$, propter quem statum determinatiōis duplex est latus, ideoq; duplex solutio.

Præter eam autem determinatiōem in eo casu laterū æqualiū, est et alia determinatio quæ statim patet ex analysi quæ inducit aliquam limitatiōem, et extra quam si æquatio propositæ nō satisfaciat latera erunt fictitia, ac proinde quæstio impossibilis, quemadmodū in his duab\textsuperscript{9} æquationib\textsuperscript{9} quadraticis $z\,\mathrm{pl.}$ nō debet esse maius \uncertain{quadrante} quadrati $x$, cū enim $x$, sit summa laterū, rectangulū maius eo non potest quam si utraq; latera sint æqualia inter se, quod quidem rectangulū eo casu æquale erit quartæ partj $x$ quadratj seu dimidio rectæ $x$, secus verò cū $x$ est differentia ut in sequentibus huius partis æquationibus.

\subsection*{Ad Caput Tertium}

In 1\textsuperscript{a} propositiōe cū latus suprà sit maius lateri infrà hæc constitutio sic Interpretabitur.
\marginal{$z\,\mathrm{pl.} + xA - A^{2}$}
\[ z\,\mathrm{pl.} + xA - A^{2} \parallel 0 \]
Et erunt duo latera, alterū idemq; maius suprà, alterū verò Idemq; minus infra, quorū differentia est $x$, rectangulū sub ipsis $z\,\mathrm{pl.}$ et $A$ fit quodlibet ex ipsis.

In 2\textsuperscript{a} verò propositiōe, cū duo latera sint Inter se æqualia evanescet affectio sub latere et huiusmodj æquatiōis hæc erit formula.
\marginal{$z\,\mathrm{pl.} - A^{2} \parallel 0$}
\[ z.\,\mathrm{pl.} - A^{2} \parallel 0 \]
Et erunt duo latera inter se æqualia alterū supra, alterū verò infrà, rectangulū sub ipsis $z\,\mathrm{pl.}$, et fit $A$, alterutrum.

In 3\textsuperscript{a} denique propositiōe cū latus infra sit maius lateri suprà, ideò hæc æquatio hanc Interpretatiōem accipiet.
\marginal{$z\,\mathrm{pl.} - xA - A^{2}$}
\[ z\,\mathrm{pl.} - xA - A^{2} \parallel 0 \]
Et sunt duo latera alterū Idemq; maius infrà, alterū Idemq; minus suprà, quorū differentia est $x$, rectangulū sub ipsis est $z\,\mathrm{pl.}$ et fit $A$ alterutrum.

\page{6}
\note{Double page ; page de gauche paginée « 6 », page de droite « 7 ».}

\subsection*{De Constitutione Æquationum Cubicarum}

Æquatio Cubica est cū potestas est cubus, ac proinde de tribus lateribus est explicabilis quæ quidem tria latera vel erunt suprà, vel duo suprà et unū Infrà, vel duo Infrà et unū suprà, vel denique tria Infrà, quæ divisio totidem scilicet quatuor dabit capita. Cum verò hæc æquatio cubica possit etiam à quadraticâ et laterali procreari, vel ex se subsistere erunt rursus et alia duo capita.

\subsection*{Caput primū In quo tria latera sunt suprà}

Propositio unica.
\[ \begin{array}{lll} \text{Sit } B - A \parallel 0 & \text{vel} & \text{vel} \\ C - A & C - A \parallel 0 & D - A \\ D - A \text{ cuiuscumq; valoris} & D - A & B - A \\ & B - A \text{ cuiuscumq; valoris} & C - A \text{ cuiuscumq; valoris} \end{array} \]
productum quod Inde oritur scilicet.
\[ BCD - CDA + DA^{2} - BDA + CA^{2} - BCA + BA^{2} - a^{3} \parallel 0 \]
\note{Les trois lignes du produit sont écrites l'une sous l'autre, $BCD$ en tête et $-a^3$ à la fin de la troisième ; elles sont ici mises bout à bout.}

Æquale nihilo, et manifestū est ex iis quæ dicta sunt tria latera esse positiva suprà, hæc â æquatio hanc accipiet Interpretatiōem.
\marginal{$T\,\mathrm{sd.} - z\,\mathrm{pl.}\,A + xA^{2} - a^{3} \parallel 0$}
\[ T\,\mathrm{Sol.} - z\,\mathrm{pl.}\,A + xA - a^{3} \parallel 0 \]
\note{Le terme en $x$ est écrit « $+\,\dot{x}A$ », sans le 2 de l'exposant, que porte la forme recopiée dans la marge ; laissé tel quel.}

Erunt Igitur tria latera suprà, quorū summa est $x$, et summa rectangulorū sub ipsis binis et binis sumptis est $z\,\mathrm{pl.}$ et solidū sub his rectangulis $T\,\mathrm{Sol.}$, et $A$, fit quodlibet ex lateribus.

\subsection*{Caput secundū In quo duo latera sunt suprà et 3\textsuperscript{ium} Infrà}

\[ \begin{array}{lll} \text{Sit } B - A \parallel 0 & \text{vel} & \text{vel.} \\ C - A & C - A \parallel 0 & D + A \parallel 0. \\ D + A & B - A & B - A \\ & D + A & C - A \end{array} \]
productum erit
\[ BCD - BDA + DA^{2} - CDA - BA^{2} + BCA - CA^{2} + A^{3} \parallel 0 \]

Quæ quidem constitutio, ob variam signorū $+$ et $-$ notatiōem, quæ præficiuntur lateribus diversæ nãtūræ scilicet iis quæ sunt suprà et alterj quod est infrà, diversam quoq; Interpretatiōem accipiet, secundū horū laterū magnitudines, siquidem vel latus infrà maius est summa laterū supra, vel ipsi æquale vel ipsi minus, et rursus ex eo varia orietur constitutio, quod tametsi latus positivū Infrà sit minus summâ laterū positivorū suprà, poterit nihilominus fieri ut rectangulum sub summa laterū suprà in latus infrà sit vel æquale, vel etiam minus rectangulo quod continetur sub lateribus suprà, et inde tot orientur propositiones.

\subsection*{Propositio 1\textsuperscript{a}}

\marginal{$T\,\mathrm{sd.} - z\,\mathrm{pl.}\,A + xA^{2} + A^{3}$}
Cum latus Infrà est maius summâ laterum suprà.

Manifestū autem est ex hac hypothesi ambo rectangula $BD + CD$ esse maiora rectangulo

\note{Page de droite, paginée « 7 ». Les formes canoniques sont recopiées dans la marge extérieure en regard de chaque proposition.}

$BC$. erit Igitur
\[ T\,\mathrm{Sd.} - z\,\mathrm{pl.}\,A + xA^{2} + A^{3} \parallel 0 \]
Et sunt tria latera duo suprà, alterum idemq; maius duobus reliquis simul sumptis infrà, $x$ est differentia horū laterum $z\,\mathrm{pl.}$ differentia rectangulorū et $T\,\mathrm{Sol.}$ quod fit sub ipsis, $A$ vero fit quodlibet ex lateribus.

\subsection*{Propos. 2\textsuperscript{a}}

\marginal{$T\,\mathrm{Sol.} - z\,\mathrm{pl.}\,A + A^{3}$}
Cum latus infra æquatur agg\textsuperscript{o} laterū supra. Patet itiam In hac propositiōe rectangula $BD + CD$ esse maiora rectangulo $BC$ erit Igitur
\[ T\,\mathrm{Sol.} - z\,\mathrm{pl.}\,A + A^{3} \parallel 0 \]
et cæt. ut supra præter quam hic nulla est affectio sub quadrato cū nulla sit differentia lateris Infrà, et summæ eorū quæ sunt supra.

\subsection*{Propos. 3\textsuperscript{a}}

\marginal{$T\,\mathrm{Sol.} - z\,\mathrm{pl.}\,A - xA^{2} + A^{3}$}
Cum latus infrà minus est summâ laterū suprà, et rectangulū sub summâ laterū suprà in latus infrà, maius est rectangulo sub lateribus suprà. Et hæc constitutio hanc Interpretatiōem accipiet.
\[ T\,\mathrm{Sol.} - z\,\mathrm{pl.}\,A - xA^{2} + A^{3} \parallel 0 \]
Et sunt tria latera duo suprà simul maiora reliquo Infrà, $x$ fit excessus horū duorū laterū, $z\,\mathrm{pl.}$ excessus rectangulorū $BD + CD$ in rectangulū $BC$, Et $T\,\mathrm{Sd.}$ est solidū quod continetur sub his lateribus et $A$, fit quodlibet ex ipsis.

\subsection*{Propos. 4\textsuperscript{a}}

\marginal{$T\,\mathrm{Sol.} - xA^{2} + A^{3}$}
Cum idem latus Infrà minus est et cæt. sed rectangula $BD + CD$ simul æqualia rectangulo $BC$ quæ constitutio sic Interpretabitur.
\[ T\,\mathrm{Sol.} - xA^{2} + A^{3} \parallel 0 \]
Et affectio sub latere evanescit, quia nulla est differentia rectangulorum.

\subsection*{Propos. 5\textsuperscript{a}}

\marginal{$T\,\mathrm{Sol.} + z\,\mathrm{pl.}\,A - xA^{2} + A^{3}$}
Cum Idem latus Infra minus est et cæt., et rectangula $BD + CD$ simul minora sunt rectangulo $BC$. Secundum quam hypothesim sequitur hoc latus Infrà minus esse quovis eorū quæ sunt suprà, ut videbimus Inferius, hæc autem constitutio sic Interpretabitur.
\[ T\,\mathrm{Sol.} + z\,\mathrm{pl.}\,A - xA^{2} + A^{3} \parallel 0 \]
Et sunt ut supra.

\subsection*{Caput Tertium}

In quo sunt duo latera Infrà et alterum suprà.
\[ \begin{array}{lll} \text{Sit } D + A \parallel 0 & \text{vel} & \text{vel} \\ \text{\& } C + A & C + A & B - A \\ B - A \text{ cuiuscumq; valoris} & B - A & C + A \\ & D + A \text{ et cæt.} & D + A \text{ Et Cætera.} \end{array} \]
Productum quod Inde oritur est
\[ BDC + BCA + BA^{2} + BDA - CA^{2} - CDA - DA^{2} - A^{3} \parallel 0 \]
\note{Le produit est écrit en trois lignes accolées après « $BDC$ » ; elles sont ici mises bout à bout.}

\page{7}
\note{Double page ; page de gauche paginée « 8 », page de droite « 9 ».}

Inde autem oriuntur quinque propositiones ut In præcedentj cap\textsuperscript{e} quæ ōes Et singulæ inversæ sunt, sic et totū caput totius inversū, quæ enim latera illj erant suprà, hic sunt infrà, et quod illic erat infrà, hic est suprà.

\subsection*{Propos. 1\textsuperscript{a}}

\marginal{$T\,\mathrm{Sol.} + z\,\mathrm{pl.}\,A + xA^{2} - A^{3}$}
Cum latus suprà maius est summâ laterū Infra quæ constitutio sic Interpretatur
\[ T\,\mathrm{Sol.} + z\,\mathrm{pl.}\,A + xA^{2} - A^{3} \parallel 0 \]

\subsection*{Propos. 2\textsuperscript{a}}

\marginal{$T\,\mathrm{Sol.} + z\,\mathrm{pl.}\,A - A^{3}$}
Cum latus suprà æquabitur agg\textsuperscript{o} laterū Infrà. hæc erit Interpretatio
\[ T\,\mathrm{Sol.} + Z\,\mathrm{pl.}\,A - A^{3} \parallel 0 \]

\subsection*{Propos. 3\textsuperscript{a}}

\marginal{$T\,\mathrm{Sol.} + z\,\mathrm{pl.}\,A - xA^{2} - A^{3}$}
Cum latus suprà minus est summâ laterū Infrà et rectangulū sub summâ laterū Infra in latus suprà maius est rectangulo sub lateribus Infrà. Adhibita Interpretatiōe erit.
\[ T\,\mathrm{Sol.} + z\,\mathrm{pl.}\,A - xA^{2} - A^{3} \parallel 0 \]

\subsection*{Propos. 4\textsuperscript{a}}

\marginal{$T\,\mathrm{Sol.} - xA^{2} - A^{3}$}
Cum Idem latus suprà minus est eâdem summâ et rectangula $BC + BD$, simul sunt æqualia rectang\textsuperscript{lo} $CD$. erit
\[ T\,\mathrm{Sol.} - xA^{2} - A^{3} \parallel 0 \]

\subsection*{Propos. 5\textsuperscript{a}}

\marginal{$T\,\mathrm{Sol.} - z\,\mathrm{pl.}\,A - xA^{2} - A^{3}$}
Cum Idem latus suprà, minus est eâdem summâ et rectangula $BC + BD$ simul minora sunt rectangulo $CD$. Cuius constitutiōis hæc erit formula.
\[ T\,\mathrm{Sol.} - Z\,\mathrm{pl.}\,A - xA^{2} - A^{3} \parallel 0 \]

\subsection*{Caput 4\textsuperscript{um}}

In quo ōia latera sunt Infra.

Hoc caput inversū est primj, et ōes affectiones notantur signo $+$ quod nunquam directe et per se, evenit. huiusmodj autem constitutiōis hæc est formula.
\marginal{$T\,\mathrm{Sol.} + z\,\mathrm{pl.}\,A + xA^{2} + A^{3}$}
\[ T\,\mathrm{Sol.} + z\,\mathrm{pl.}\,A + xA^{2} + A^{3} \parallel 0 \]
Quæ quidem formula eodem modo enuntiabitur quo In cap\textsuperscript{e} 1\textsuperscript{o}. modo ea dicantur hic de lateribus infrà, quæ ibj dicta sunt de lateribus suprà.

\subsection*{Caput 5\textsuperscript{um}. De constitutiōe Æquationū Cubicarū quæ oriuntur Ex duabus Æquationibus alterâ planâ, alterâ lateralj}

Æquatiōes quæ oriuntur ex æquatiōe quadraticâ, de duobus lateribus explicabilj, recidunt In præcedentes constitutiones cubicas, cū enim æquatio quadratica ex hypothesj duas habeat radices, Intelligantur esse $B$ et $C$, deinde verò ad efformandam cubicam, de qua in hoc Capite, assumatur lateralis, cuius radix sit $D$, igitur quæ inde oritur æquatio erit explicabilis de tribus lateribus, ac proinde eius constitutio reperietur Inter eas de quib\textsuperscript{9} egimus superioribus Capitib\textsuperscript{9} huius partis.

Quia verò unica est \struck{tantū} æquatio quadratica de unico tantū latere explicabilis scilicet $B\,\mathrm{pl.} + A^{2} \parallel 0$, in quâ æquatiōe latus est positivū Infrà, si ad

\note{Page de droite, paginée « 9 ». Un chiffre « 9 » de la série courante dans la marge de gauche, qui est ici la marge intérieure.}

efformandam cubicam assumatur lateralis, hæc tantū duplex esse potest vel $C - A \parallel 0$, quo casu nempe latus est positivū supra vel $C + A \parallel 0$ ubj latus est Infrà, unde et duæ tantū huius capitis propositiones.

\subsection*{Prop. 1\textsuperscript{a}}

\marginal{$T\,\mathrm{Sol.} - Z\,\mathrm{pl.}\,A + xA^{2} - A^{3}$}
Cum æquatio quadratica $B\,\mathrm{pl.} + A^{2} \parallel 0$ ducitur In lateralem, cuius latus est positivū suprà.
\[ \begin{array}{ll} \text{Sit } B\,\mathrm{pl.} + A^{2} \parallel 0 & \text{vel. Sit } C - A \parallel 0 \\ \text{\& } C - A \text{ cuiuscumq; valoris} & \text{Et } B\,\mathrm{pl.} + A^{2} \text{ cuiuscumq; valoris} \end{array} \]
productū quod Inde oritur erit.
\[ B\,\mathrm{pl.}\,C - B\,\mathrm{pl.}\,A + CA^{2} - A^{3} \parallel 0 \]
Quæ quidem constitutio, adhibitâ scilicet convenientiore in superioribus Interpretatiōe sic erit.
\[ T\,\mathrm{Sol.} - Z\,\mathrm{pl.}\,A + xA^{2} - A^{3} \parallel 0 \]
Cuius constitutiōis formula, quamvis ōo similis sit formulæ quæ in 1\textsuperscript{o} capite huius partis, non tamen similiter enuntiabitur, $x$, enim ut videre est ex eius origine non est summa laterū, nec $z\,\mathrm{pl.}$ aggregatū rectangulorū et cæt., ac proinde hæc constitutio est exceptio regulæ generalis ex quâ certū est in ōibus æquationibus alternatim affirmatis et negatis à principio usq; ad finem, coefficientē altioris gradus sub potestate fieri summam laterū, coefficientē verò gradus Immediate subsequentis summam rectangulorū vel planorū sub iisdem lateribus quæcumq; vera ea latera ducta fuerint, similiter coefficientē gradus proximè Inferioris fieri summam omniū solidorū sub iisdem lateribus, et sic in Infinitum.

Sunt Igitur in hac formula duo latera alterū suprà æquale $x$, alterū infrà potentiâ æquale $z\,\mathrm{pl.}$ et $T\,\mathrm{Sol.}$ fit ex quadrato eius quod est infrà in altitudinem illius quod est suprà, et $A$ est quodlibet ex lateribus.

\subsection*{Propos. 2\textsuperscript{a}}

\marginal{$T\,\mathrm{Sol.} + z\,\mathrm{pl.}\,A + xA^{2} + A^{3}$}
In qua sunt duo latera ambo Infrà.
\[ \begin{array}{ll} \text{Sit } B\,\mathrm{pl.} + A^{2} \parallel 0 & \text{vel } C + A \parallel 0 \\ \text{\& } C + A \text{ cuiuscumq; valoris} & B\,\mathrm{pl.} + A^{2} \text{ cuiuscumq; valoris} \end{array} \]
productū erit
\[ B\,\mathrm{pl.}\,C + B\,\mathrm{pl.}\,A + CA^{2} + A^{3} \parallel 0 \]
Cuius constitutiōis hæc esto Interpretatio.
\[ T\,\mathrm{Sol.} + z\,\mathrm{pl.}\,A + xA^{2} + A^{3} \parallel 0 \]
Quæ etiam formula ōo similis est quoad signa formulæ quartj capitis huius partis, sed maximè Inter se differunt, ut patet ex diversâ earū origine, igitur huius enunciatio eadem erit quæ in præcedentj nisi quod hic ambo latera sunt Infra.

\page{8}
\note{Double page ; page de gauche paginée « 10 », page de droite « 11 ».}

\subsection*{Caput 6.\textsuperscript{um} De Æquationibus Cubicis quæ per se subsistunt}

Cum In præcedentibus capitibus mixta sint omnimodè tria latera, deinde etiam omnimode planū et latus ad constituendas æquationes cubicas, si præter hanc lateralem veriam vel mixtiōem vel Combinatiōem et æquationes quæ inde ortæ sunt quædam efficiantur, manifestū est illas nec ex tribus nec ex duabus æquationibus constare, atq; ità neque de tribus, neq; de duobus lateribus explicabiles, esse.

Reperiuntur eiusmodj æquatiōes omnino sex unde totidem propositiones constituentur.

\marginal{$T\,\mathrm{Sol.} + xA^{2} + A^{3}$}
\subsection*{Prop. 1\textsuperscript{a}}
\[ T\,\mathrm{Sol.} + xA^{2} + A^{3} \parallel 0 \]
\marginal{$T\,\mathrm{Sol.} + z\,\mathrm{pl.}\,A + A^{3}$}
\subsection*{Prop. 2\textsuperscript{a}}
\[ T\,\mathrm{Sol.} + z\,\mathrm{pl.}\,A + A^{3} \parallel 0 \]
\marginal{$T\,\mathrm{Sol.} + A^{3}$}
\subsection*{Prop. 3\textsuperscript{a}}
\[ T\,\mathrm{Sol.} + a^{3} \parallel 0 \]
In his tribus propositionibus est unicū latus Infrà.

\marginal{$T\,\mathrm{Sol.} - z\,\mathrm{pl.}\,A - A^{3}$}
\subsection*{Prop. 4\textsuperscript{a}}
\[ T\,\mathrm{Sol.} - z\,\mathrm{pl.}\,A - A^{3} \parallel 0 \]
\marginal{$T\,\mathrm{Sol.} - xA^{2} - A^{3}$}
\subsection*{Prop. 5\textsuperscript{a}}
\[ T\,\mathrm{Sol.} - xA^{2} - A^{3} \parallel 0 \]
\subsection*{Prop. 6\textsuperscript{a}}
\[ T\,\mathrm{Sol.} - A^{3} \parallel 0 \]
In his tribus ultimis propositionibus est unicum latus suprà.

\subsection*{De Determinationibus Æquationum Cubicarum}

Præmissis constitutionibus æquationū cubicarū reliquū est ut de earum determinationib\textsuperscript{9} agamus. Cū dictæ determinatio duplex sit ut notavimus, maior scilicet et minor, \& maior sit cū omnia latera sunt æqualia, minor verò cum quædam tantū, manifestū erit ex harū æquationū constitutionib\textsuperscript{9} maiorj determinationj solummodo esse obnoxias formulas primj et quartj capitis, quæ etiam erunt obnoxiæ et minorj, cū fierj possit ut in his æquationibus non omnia latera sint æqualia sed quædam duntaxat.

Quoniam verò præter has determinatiōes aliam notavim\textsuperscript{9} quæ limitatiōem inducit, et extra quæstio proposita resolvj nō poterit, lateraq; erunt merè fictitia, ideò In constitutionibus duorū dictorū capitū in quibus $x$ est summa laterum eo casu quæstio erit possibilis si $Z\,\mathrm{pl.}$ non sit maius tertia parte $x$ quadratj, et similiter $T\,\mathrm{pl.}$ non sit maius tertiâ parte cubj eiusdem $x$, si enim linea in tres partes dividatur maius solidū sub his tribus portionibus contentum, erit cū dictæ

\note{Page de droite, paginée « 11 ».}

tres partes erunt Inter se æquales, Similiter eâdem lineâ adhuc trifariam sectâ maior summa rectangulorū seu planorū, iis binis et binis sumptis erit cū eâdem eiusdem lineæ partes erunt inter se æquales, quod ita evidens est, ut maiorj demonstratiōe non egeat.

Maior itaq; determinatio primj et quartj capitis erit cū $x$ Inter partes æquales dividetur.

Si verò In his capitibus sicut et in sequentib\textsuperscript{9}, tria latera non sint inter se æqualia, sed duo tantum, hæc duo latera sic investigabuntur.

Sint enim In constitutiōe primj capitis $B$ et $C$, hæc duo latera æqualia patet productū huius constitutiōis ad hoc reducj.
\[ B^{2}D - 2BDA + 2BA^{2} - B^{2}A + DA^{2} - A^{3} \parallel 0 \]
\note{Le produit est écrit en deux lignes ; elles sont ici mises bout à bout.}

Et similj adhibitâ interpretatiōe $T\,\mathrm{Sol.}$ fierj $B^{2}D$, et tria plana $2BD$ et $B^{2}$ similiter fierj $Z\,\mathrm{pl.}$ et $x$ esse summam triū laterū $2B + D$, quorū duo erunt ex hypothesi Inter se æqualia, quæ ut dignoscantur. alterū eorū sit $A$, igitur $2A$ erunt duo latera æqualia $B$ et $C$, \& $x - 2A$ erit alterū latus prioribus inæquale, ut autem sub speciebus eorumdem laterū æqualiū habeatur $Z\,\mathrm{pl.}$ quoniam innotescit ex iam proposita constitutiōe idem $Z\,\mathrm{pl.}$ æquari bis rectangulo sub latere æqualj in illud quod est inæquale $+$ quadrato eius quod est æquale, si igitur $x - 2A$ ducatur in $2A$ et producto adiiciatur $A^{2}$ erit
\[ 2xA - 3A^{2} \parallel Z\,\mathrm{pl.} \]
Quorū omniū si sumatur tertia pars erit
\[ \tfrac{2}{3}xA - A^{2} \parallel \tfrac{1}{3}Z\,\mathrm{pl.} \quad \text{atq; ideo} \]
\[ \tfrac{1}{3}Z\,\mathrm{pl.} - \tfrac{2}{3}xA + A^{2} \parallel 0 \]
Et sic cū sese offerat æquatio cognita de duob\textsuperscript{9} lateribus explicabilis, de qua in recognitiōe quadraticarū factū est satis propositæ disquisitionj, cui $A$ enuntiatur de alterutro latere, \& inductio legitima ut notavimus cū numerus partiū est finitus, et si $T\,\mathrm{Sol.}$ esset ignotū nihil amplius agendū superesset, hoc et $A$, habeatur unicū $T\,\mathrm{Sol.}$ ad hoc ut \ill{} \uncertain{proprietas} laterū quæ investigantur reducatur, quod quidem $T\,\mathrm{Sol.}$ cū ex constitutiōe proposita æquetur quadrato lateris æqualis in latus inæquale, erit
\[ x\,A^{2} - 2A^{3} \parallel T\,\mathrm{Sol.} \quad \text{et sumptis dimidiis} \]
\[ \tfrac{1}{2}xA^{2} - A^{3} \parallel \tfrac{1}{2}T\,\mathrm{Sol.} \quad \text{et per Antithesim} \]
\[ \tfrac{1}{2}T\,\mathrm{Sol} - \tfrac{1}{2}xA^{2} + A^{3} \parallel 0 \]
Quæ æquatio cū æquetur nihilo, sicut et ea quam habuimus de $Z\,\mathrm{pl.}$ et $A$ sit Idem in utrâq;, Igitur ambæ hæ duæ æquatiōes poterunt Inter se conferrj ut habeatur ipsū $A$, ac proinde ut sint homogeneæ.
\[ \tfrac{1}{3}z\,\mathrm{pl.}\,A - \tfrac{2}{3}xA^{2} + A^{3} \parallel \tfrac{1}{2}T\,\mathrm{Sol.} - \tfrac{1}{2}xA^{2} + A^{3} \]
et ablato utrimq; $A^{3}$ erit
\[ \tfrac{1}{3}z\,\mathrm{pl.}\,A - \tfrac{2}{3}xA^{2} \parallel \tfrac{1}{2}T\,\mathrm{Sol} - \tfrac{1}{2}xA^{2} \quad \text{et per Antithesim} \]
\[ \tfrac{1}{2}T\,\mathrm{Sol.} - \tfrac{1}{3}z\,\mathrm{pl.}\,A + \tfrac{1}{6}xA^{2} \parallel 0 \]
quorum omnium si sumatur sexta pars et productū applicetur ad $x$ erit
\[ \frac{3\,T\,\mathrm{Sol}}{x} - \frac{2\,z\,\mathrm{pl.}\,A}{x} + A^{2} \parallel 0 \]

\page{9}
\note{Double page ; page de gauche paginée « 12 », page de droite « 13 ». Un chiffre « 10. » de la série courante dans la marge extérieure de la page de gauche.}

Cuius æquatiōis rursus alia Instituenda est comparatio cū dictâ 1\textsuperscript{a} quadratica de $z\,\mathrm{pl.}$ ut tandem habeatur $A$, unicū, ambæ enim sicut et supra æquantur nihilo, habebimus Igitur.
\[ \frac{3\,T\,\mathrm{Sol}}{x} - \frac{2\,Z\,\mathrm{pl.}\,A}{x} + A^{2} \parallel \tfrac{1}{3}z\,\mathrm{pl.} - \tfrac{2}{3}x\,A + A^{2} \]
utrimque auferatur $A^{2}$. tum erit
\[ \frac{3\,T\,\mathrm{Sol}}{x} - \frac{2\,Z\,\mathrm{pl.}\,A}{x} \parallel \tfrac{1}{3}Z\,\mathrm{pl.} - \tfrac{2}{3}xA \quad \text{Et per Antithesim.} \]
\[ \frac{3\,T\,\mathrm{Sol}}{x} - \tfrac{1}{3}Z\,\mathrm{pl.} \parallel \frac{2\,Z\,\mathrm{pl.}\,A}{x} - \tfrac{2}{3}xA \quad \text{Ac proinde} \]
\[ \frac{\dfrac{3\,T\,\mathrm{Sol}}{x} - \tfrac{1}{3}Z\,\mathrm{pl.}}{\dfrac{2\,Z\,\mathrm{pl.}}{x} - \tfrac{2}{3}x} \parallel A \]

Oportet Igitur in casu determinatiōis ut alterū ex duobus lateribus æqualib\textsuperscript{9} oriatur ex hac applicatiōe, et sic in ōib\textsuperscript{9} constitutionib\textsuperscript{9} sive cubicis sive quadratoquadratis, sive altioris gr\textsuperscript{is}, eodem plane artificio latus æquale detegetur, nec maior erit dificultas, nisi quod in iis constitutionib\textsuperscript{9} ubj potestas fuerit in altiori gradu plures instituenda erunt æquationes, ut tandem ad ultimam deveniatur in quâ ut in proposito exemplo habeatur $A$, unicū.

His ita notatis quæ spectant in universū ad determinationes non solū harū æquationū cubicarū sed etiam aliarū sive quadratoquadraticarū sive etiam gradus altioris, sigillatim videndū est de determinationibus, in casu laterū æqualiū quæ ad unamquamq; propositiōem ōniū capitū huius partis pertinent, et etiam de iis determinationibus, quæ limitatiōem inducunt.

Et in 1\textsuperscript{o} quidem dictū est de iis ōib\textsuperscript{9} quæ spectant tū ad Caput 1\textsuperscript{um} tū ad caput quartū. Consequens est ut agamus de determinationib\textsuperscript{9} unicuiq; propositiōj capitis secundj inservientib\textsuperscript{9}. Et statim patet has propositiones sicut et reliquas capitis tertij, quintj et sextj maiorj quam diximus determinationj nō esse obnoxias, in iis enim constitutionib\textsuperscript{9} vel unū latus maius est duob\textsuperscript{9} reliquis simul, vel minus, vel ipsis simul sumptis æquale, vel etiam duo tantū sunt latera, vel certè unicū ut In constitutionib\textsuperscript{9} capitis ultimj. Videndū est Igitur de determinatiōe minorj, et de ea quæ limitatiōem aliquam inducit.

Prop\textsuperscript{is} 1\textsuperscript{a}. Capitis 2.\textsuperscript{i} determinatio quæ limitatiōem inducat nulla est, patet enim ex analysj et ex applicatiōis planj per excessū, quæ quidem applicatio semper est possibilis, determinatio vero in casu duorū laterū æqualiū orietur ex hac applicatiōe
\[ \frac{\dfrac{3\,T\,\mathrm{Sol}}{x} + \tfrac{1}{3}Z\,\mathrm{pl.}}{\dfrac{2\,Z\,\mathrm{pl.}}{x} + \tfrac{2}{3}x} \]
Et $A$, alterutrū ex duobus lateribus suprà æquabitur longitudinj quæ Inde proveniet.

Prop\textsuperscript{is} 2\textsuperscript{a} Cap\textsuperscript{is} 2.\textsuperscript{i} determinatio sic elicietur. Quoniam enim In hac constitutiōe duo latera suprà simul sumpta sunt æqualia lateri infra, ideoq; affectio sub

\note{Page de droite, paginée « 13 » à l'encre en haut à l'angle extérieur. Les chiffres « 11 », « 12 » et « 13 » de la série courante sont dans la marge de gauche, qui est ici la marge intérieure ; le dernier porte une note en français, seule note non latine de ces vingt vues.}

quadrato evanescit, si statuamus hoc latus infra scilicet $D$, æquari $2A$, Igitur $B$, erit $A$, et $C$ similiter erit $A$, ac proinde $z\,\mathrm{pl.}$ quod diximus esse differentiam planorū $BD + CD$, et $BC$, sive excessū quo planū sub $B + C$ in $2D$, maius est plano $BC$, erit æquale $3A^{2}$. et ideo
\[ A^{2} \parallel \tfrac{1}{3}\,z\,\mathrm{pl.} \]
secundū quam Interpretatiōem $z\,\mathrm{pl.}$ fit summa triū planorū proportionaliū; scilicet $B^{2} + BC$ et $C^{2}$ quorū ōiū idem latus scilicet $A$.

vel etiam Idem $z\,\mathrm{pl.}$, hanc Interpretatiōem accipiet ut sit differentia quadrati lateris Infrà et rectanguli lateriū suprà.

Similiter et $T\,\mathrm{Sol.}$ tribus etiam variis modis poterit enuntiari. hoc præmisso Lemmate.

\subsection*{Lemma}

Si sint duo latera $B$ et $C$ à quibus fiant tria plana proportionalia, $B^{2}$, $BC$ et $C^{2}$, tria solida quæ ab his planis efficientur, erunt æqualia, primū quod à plano Intermedio et summā lateriū extremorū, secundū quod à duobus primis in altitudinem lateris eius quod est extremū, tertiū denique quod à duobus extremis in altitudinem lateris primi.

Cū enim hæc tria plana sint proportionalia erit ut $B^{2}$ ad $BC$ sic $B$ ad $C$ et componendo ut $B^{2} + BC$ ad $BC$. sic $B + C$ ad $C$., igitur constat \uncertain{ꝓportū} ut solidum cuius basis $B^{2} + BC$ et altitudo $C$, æquale solido cuius basis $BC$ altitudo vero $B + C$, eademq; ratiōe huic erit æquale solidū cuius basis $C^{2} + BC$ in altitudinem $B$, ac proinde tria hæc solida Inter se æqualia, et $T\,\mathrm{Sol.}$ erit quodlibet ex ipsis, fit enim in casu propositæ determinatiōis ex $B$ et $C$ Inter se æqualibus, et $D$, quod etiam ambobus simul sumptis est æquale.
\note{Vérification numérique du lemme, écrite dans la marge de gauche d'une plume plus fine, sur $b = 4$ et $c = 10$ ; la lettre $c$ y est tracée comme un $r$ ouvert, et un nombre raturé reste illisible sous la biffure, au-dessous du premier « 560 ». Elle est transcrite ci-après.}
\[ \begin{array}{lll} b & & c \\ 4 & & 10 \\ \hline b^{2} & bc & c^{2} \\ 16. & 40 & 100 \\ \hline b^{2}\ \ bc & \parallel & b\,.\,c \\ 16\ \ 40 & & 4\,.\,10 \\ \hline b^{2}+bc\ \ bc & \parallel & b+c \quad c \\ 56 \qquad 40 & & 14 \qquad 10 \end{array} \]
\[ \begin{array}{r} 56 \\ 10 \\ \hline 560 \end{array}\ \begin{array}{l} b^{2}+bc \\ c \end{array} \qquad \begin{array}{r} 14 \\ 40 \\ \hline 560 \end{array}\ \begin{array}{l} b+c \\ bc \end{array} \qquad \begin{array}{r} 140 \\ 4 \\ \hline 560. \end{array}\ \begin{array}{l} c^{2}+BC \\ B \end{array} \]

In hac Igitur constitutiōe ea est determinatio ut latus $T\,\mathrm{Sol.}$ æquetur \add{potentia} $\tfrac{1}{3}\,z\,\mathrm{pl.}$ et hæc est limitatio \uncertain{ꝑ} $T\,\mathrm{Sol.}$ non fit maius solido cuius basis $\tfrac{2}{3}\,z\,\mathrm{pl.}$ altitudo vero est $\tfrac{1}{3}$ eiusdem planj quæ quidem determinatio patet ex analysi. eximemus ex hac hypothesi $T\,\mathrm{Sol.} \parallel 2A^{3}$. quod quidem solidū æquatur $z\,\mathrm{pl.}\,A - A^{3}$, Igitur facta antithesi $3A^{3} \parallel z\,\mathrm{pl.}\,A$. et ōibus aplicatis ad $A$, $3A^{2} \parallel z\,\mathrm{pl.}$ et ex consequenti $A^{2}$ potentia scilicet lateris $T\,\mathrm{Sol.}$ æquabitur $\tfrac{1}{3}\,z\,\mathrm{pl.}$. Si deinde prioris æquationis in qua $T\,\mathrm{Sol.} \parallel 2A^{3}$. dimidium ex utraq; parte sumatur \struck{erit}
\note{« potentia » est ajouté au-dessus de la ligne, appelé par un caret.}
\[ \text{erit } \tfrac{1}{2}T\,\mathrm{Sol.} \parallel A^{3} \quad \text{et ōibus aplicatis ad } A \]
\[ \text{erit } \frac{\tfrac{1}{2}T\,\mathrm{Sol.}}{A} \parallel A^{2} \quad \text{atq; ideo} \]
\[ \frac{T\,\mathrm{Sol.}}{A} \parallel 2A^{2} \quad \text{et ex consequenti} \]
\[ \frac{T\,\mathrm{Sol.}}{A} \parallel \tfrac{2}{3}\,z\,\mathrm{pl.} \quad \text{ōibus ductis In } A. \]
\[ \text{erit } T\,\mathrm{Sol.} \parallel \tfrac{2}{3}\,z\,\mathrm{pl.}\,A \quad \text{et manifestū est } A \text{ fieri } \sqrt{\tfrac{1}{3}\,z\,\mathrm{pl.}} \]

Sed eadem determinatio aliter sic demonstrabitur præmisso Lemmate quod a D. Rob. demonstratū est., Et hoc Lemma videbitur pagina sequentj.
\marginal{13 en son traité de mechanique des plans inclinez}
\note{La note marginale est en français ; elle renvoie, comme « D. Rob. » de la ligne qu'elle accompagne, à Roberval et à son traité de mécanique des plans inclinés. Le bas de la page est resté blanc.}

\page{10}
\note{Double page ; page de gauche paginée « 14 », page de droite « 15 ». Les chiffres « 14 », « 15 », « 16 » et « 17 » de la série courante sont dans la marge de gauche.}

\subsection*{Lemma}

Omniū parallelepipedorū secundū æqualia plana, vel Idem planū parallelogrammū aplicatorū deficientiumq; figuris parallelepipedis similibus Inter se et similiter positis, maximū est quod ad duas tertias aplicatur duplex existens defectus.

Ex quo Lemmate hoc deducitur quasi corollariū.

Omne planū duabus rectis lineis inæqualibus continetur, quarū minor potens est tertiæ partis eiusdem planj, maior vero sit tripla minoris.

Quod quidem verū est In qualibet aliā \uncertain{ꝑte}, seu ut facilius explicetur in qualibet aliā fractiōe numericà, si referatur ad denominatorem ita ut maior linea sit æquemultiplex minoris, ac ipsa minor \struck{potest} potens est partis eius planj cuius ipsamet est minus latus, si enim v.g. minor possit quartam cuiusdam planj, maior erit quadrupla ipsius minoris.

His ut demonstratis ōe planū atq; ideo $z\,\mathrm{pl.}$, sub duabus rectis continetur quarū minor est potentiâ æqualis tertiæ p\textsuperscript{tj} eiusdem planj, maior vero est tripla minoris, rectangulū igitur sub eâdem minore, et duplo ipsius æquabitur duabus tertiis $z\,\mathrm{pl.}$ sumptaq; cōi altitudine eâdem minori, solidū cuius basis $\tfrac{2}{3}\,z\,\mathrm{pl.}$ altitudo vero est $\tfrac{1}{3}$ eiusdem $z\,\mathrm{pl.}$, æquabitur duabus tertiis solidj cuius basis idem $z\,\mathrm{pl.}$, altitudo vero est $\tfrac{1}{3}$ eiusdem, ac proinde si huiusmodj solidū aplicetur ad duas tertias dicti $z\,\mathrm{pl.}$, nec excedet nec deficiet, sed si ad idem $z\,\mathrm{pl.}$ aplicetur deficiet figurâ cubicâ cuius basis $\tfrac{1}{3}\,z\,\mathrm{pl.}$, altitudo vero est $\tfrac{1}{3}\,z\,\mathrm{pl.}$, duplex existens defectus, sed omniū parallelepipedorum, secundū æqualia plana et cæt. maximū est quod ad duas tertias ipsius aplicatur duplum existens defectus, Igitur solidū cuius basis $\tfrac{2}{3}\,z\,\mathrm{pl.}$ altitudo vero est $\tfrac{1}{3}$ eiusdem planj, maximū est ōiū solidorū secundū $z\,\mathrm{pl.}$ aplicatorū, deficientiumq; figurâ cubicâ, sed $z\,\mathrm{pl.}\,A - A^{3}$, est solidū aplicatū secundū $z\,\mathrm{pl.}$, et deficit figurâ cubicâ; Igitur vel $z\,\mathrm{pl.}\,A - A^{3}$, minus est solido cuius basis $\tfrac{2}{3}\,z\,\mathrm{pl.}$ altitudo vero est $\tfrac{1}{3}$ eiusdem planj, vel si ipsi æquale maximū est ōiū secundū $z\,\mathrm{pl.}$ aplicatorū, deficientiumq; figurâ cubicâ, sed $T\,\mathrm{Sol.}$ facta antithesi æquatur $z\,\mathrm{pl.}\,A - A^{3}$. Igitur vel $T\,\mathrm{Sol.}$ est æquale solido cuius basis $\tfrac{2}{3}\,z\,\mathrm{pl.}$ altitudo vero est $\tfrac{1}{3}$ eiusdem planj, vel certè ipso minus quod est propositū.

Propositionis 3\textsuperscript{æ} Cap. 2.\textsuperscript{i} determinatio potest esse multiplex, cū enim In hac constitutiōe sint tria latera duo suprà, alterū infrà, idemq; minus duobus suprà simul sumptis, vel illud latus infrà erit æquale unicuiq; eorū quæ sunt suprà, vel alterutrj tantū, vel etiam neutrj sed hæc duo latera suprà erunt Inter se æqualia.

Primo casu erit locus determinationj maiorj et $T\,\mathrm{Sol.}$ æquabitur $x^{3}$ sicut et $z\,\mathrm{pl.}$ quadrato ipsius $x$.

Secundo casu cū hoc latus infrà, scilicet $D$, sit æquale vel $B$, vel $C$, si statuamus v.g. $B$, esse $A$, erit etiam $D$ ipsi $A$, quia vero latera suprà, in hac constitutione, simul sumpta sunt maiora eo latere infrà, quorū differentia est $x$, consequens est ut $C$, alterū latus suprà æquetur $x$, et $z\,\mathrm{pl.}$ similiter æquetur quadrato ex $A$, cū rectangulū $BD$ sit differentia rectangulorū $BD + CD$, et $BC$, Ideoq; et \struck{$T\,\mathrm{Sol.}$}
\[ T\,\mathrm{Sol.} \parallel A^{2}x \quad \text{quo aplicato ad } x\text{, erit} \quad \frac{T\,\mathrm{Sol.}}{x} \parallel A^{2} \ \text{sive } z\,\mathrm{pl.} \]
hinc Igitur in dicto hoc secundo casu determinatio ut ex aplicatiōe $T\,\mathrm{Sol.}$ ad $x$, oriatur $z\,\mathrm{pl.}$

Tertio denique casu cū latera suprà erunt Inter se æqualia, eâdem utemur methodo, quâ in superioribus, et tandem post aliquot æquationes alteram ex

\note{Page de droite, paginée « 15 ». Le chiffre « 18 » de la série courante dans la marge de gauche.}

Lateribus æqualibus orietur ex hac aplicatiōe
\[ \frac{\dfrac{3\,T\,\mathrm{Sol.}}{x} - \tfrac{1}{3}\,z\,\mathrm{pl.}}{\dfrac{2\,z\,\mathrm{pl.}}{x} + \tfrac{2}{3}x} \]

Propositionis 4\textsuperscript{æ} Cap. 2.\textsuperscript{i} determinatio est ut latus infrà, minus summâ lateriū suprà, non sit maius quartâ p\textsuperscript{te} eiusdem summæ, cū enim In hac constitutiōe rectangulū sub summâ laterū suprà in latus infrà debeat æquari rectangulo sub dictis duobus lateribus suprà, manifestū est rectangulū $BC$, seu rectangulū sub lateribus suprà, eo casu esse maximū quo $B$ et $C$ erunt inter se æqualia, quo casu rectangulū $BC$, æquabitur quartæ partj quadratj rectæ ipsius $B + C$, sed quarta pars quadratj rectæ $B + C$, æquatur rectangulo sub totâ $B + C$, et quartâ p\textsuperscript{te} eiusdem, Igitur rectangulū $BC$, eo casu erit maximū, quo recta $B + C$ ducitur in quartam ipsius p\textsuperscript{tem}, Igitur $D$, per quod hæc linea $B + C$ ducitur, et ex quâ multiplicatiōe oritur rectangulū æquale rectangulo $BC$, nō debet esse maius quartâ p\textsuperscript{te} eiusdem rectæ $B + C$, hoc est quartâ p\textsuperscript{te} eiusdem summæ.

Si igitur in hac constitutiōe latera suprà sint inter se æqualia $A$, unum ex his duobus Lateribus æquabitur $\tfrac{2}{3}x$, cū enim hæc duo latera suprà sint inter se æqualia, et utrumq; sit $A$, et $D$ non possit esse maius quartâ p\textsuperscript{te} summæ duorū horū laterū scilicet $2A$, erit $D$, æquale $\tfrac{1}{2}A$, cum autem $x$, sit differentia inter $2A$ \struck{et} $\tfrac{1}{2}A$, sequitur ipsum $x$, æquari $\tfrac{3}{2}A$, ac proinde $\tfrac{2}{3}x \parallel A$.

vel certè $x$ itâ dividatur in duas p\textsuperscript{tes} inæquales ut una pars sit dupla alterius, quo casu solidū quod fit ex quadrato maioris in minorem non erit minus $T\,\mathrm{Sol.}$ in casu scilicet huius constitutionis, quod demonstrabitur præmissis duobus Lemmatibus etiam à D. Rob. demonstratis.

\subsection*{Lemma primum}

Si recta linea in duas partes inæquales dividatur, maius solidum quod fieri potest ex quadrato unius In altitudinem alterius, est cū portio eius, cuius quadratū est basis dicti solidj, dupla efficitur alterius portionis.

\subsection*{Lemma secundum}

Omnium parallelepipedorū secundū \struck{æquales rectas} eandem lineam rectam, vel secundū æquales rectas aplicatorum, deficientiumq; figuris parallelepipedis similibus Inter se et similiter positis, maximū est quod ad tertiam partem aplicatur dimidium existens defectus.

His positis demonstramus $T\,\mathrm{Sol.}$ in hac constitutione nō esse maius solido, quod fit, $x$, In duas p\textsuperscript{tes} inæquales diviso, quarū altera, alterius sit dupla, ex quadrato maioris dictarū p\textsuperscript{tiū} in altitudinem minoris.

\struck{Quoniam} Quoniam enim solidū de quo quæritur, sive solidū $CAD$ fit ex maiori portione bis sumptâ in minorem semel, in eo lateribus eius solidj scilicet $AD$, sit maior portio, scilicet duæ tertiæ partes lineæ $AB$; vel $x$, et minor portio sit æqualis $DB$, seu tertiæ parti $x$, si eiusmodj solidū aplicetur ad $AC$ sane nec excedet, nec deficiet, recta enim $AC$, est unum ex lateribus suæ basis, sed si aplicetur ipsi $AB$, vel $x$, deficiet solido cuius quodlibet latus baseos, et altitudo erunt recta $AD$, seu $CB$, sive etiam duæ tertiæ p\textsuperscript{tes} ipsius $x$, deficiet Igitur figurâ cubicâ, sed idem
\marginal{N\textsuperscript{a}. quod linea $DB$ debet esse æqualis $CD$. sed defectu cartæ non potuit ulterius produci}
\note{La figure occupe l'angle extérieur inférieur : un parallélépipède tracé en perspective, dont la base repose sur une droite lettrée, de gauche à droite, $A$, $C$, $D$, $B$ ; les sommets visibles portent $D$ (deux fois, en haut et à mi-hauteur, du côté gauche), $A$ et $C$. La note marginale ci-dessus, écrite au pli, prévient que la figure n'a pu être poussée plus loin faute de papier.}

\page{11}
\note{Double page ; page de gauche paginée « 16 », page de droite « 17 ». Le chiffre « 19. » de la série courante dans la marge de gauche.}

solidū est dimidiū defectus, quia ex tribus lineis sub quibus hoc solidū deficiens continetur, unū ex lateribus baseos et altitudo sunt eadem linea ac defectus, tertia verò sub quâ idem solidū deficiens etiam continetur est duplū eius sub quâ hoc solidū continetur, Igitur ex lemmate 2\textsuperscript{o}, hoc solidū est ōiū maximū secundū rectam $AB$, seu $x$ aplicatorum, deficientiumq; figurâ cubicâ; sed ex lemmate 1\textsuperscript{o} est etiam maximū solidū omnium quæ, ($x$ In duas p\textsuperscript{tes} inæquales diviso) continentur sub quadrato maioris p\textsuperscript{tis} in altitudinem alterius, Igitur quodcumq; solidū ad $x$ aplicetur, deficiens figurâ cubicâ, vel æquabitur eo quod fit ex quadrato maioris p\textsuperscript{tis} ipsius $x$, in minorem, cū scilicet maior pars est dupla minoris, vel certè est minus: sed $xA^{2} - A^{3}$, est solidū quod aplicatur ad $x$, deficiensq; figurâ cubicâ, ipsi $A^{3}$, æquali. Igitur ex consequenti vel solidum $xA^{2} - A^{3}$ est æquale solido cuius basis duæ tertiæ ipsius $x$, altitudo vero tertia pars eiusdem, vel ipso minus sed $xA^{2} - A^{3}$, æquatur $T\,\mathrm{Sol.}$, Igitur $T\,\mathrm{Sol.}$ vel est æquale eidem solido, scilicet $\tfrac{4}{27}x^{3}$. vel ipso minus, quod erat propositū.

Quæ quidem determinatio potest et aliter demonstrari, nec opus erit 2\textsuperscript{o} præmissorū lemmatū.

Sint enim tria latera $B$, $C$, $D$, secundū hypothesim huius constitutionis, et excessus $B + C$, in $D$ sit $x$, et rectangulū quod continetur sub summâ duorū priorum laterū in tertiū, æquetur rectangulo quod fit ex duobus primis. Dico solidū quod fit ex his tribus lateribus, scilicet $T\,\mathrm{Sol.}$ non esse maius solido cuius basis quadratum duarum tertiarū ipsius $x$, altitudo verò tertia pars eiusdem.

Nam latus $B$, æquetur rectæ $AB$, et latus $C$, rectæ $BC$, ita ut tota $AC$, sit $B + C$, latus verò $D$, sit recta $CD$, et rectangulū $ACD$, æquetur rectangulo $ABC$, vel $AF$, Igitur vel recta $CD$ æqualis erit quartæ partj rectæ $AC$, vel ipsâ minor, maior enim esse non potest ex demonstratis.
\note{Deux figures occupent la moitié extérieure de la page, côte à côte : deux rectangles dressés, le premier titré « 1\textsuperscript{us} casus », le second « 2\textsuperscript{us} casus ». Le sommet porte $A$ ; sur le côté gauche $F$, sur le côté droit $B$, puis $E$ ; la seconde figure porte en outre $G$ entre $B$ et $E$. La base est lettrée $C$ et $D$.}

Sit primo recta $CD$, æqualis quartæ partj rectæ \uncertain{$AE$}, igitur et $AF$, est quarta pars quadratj rectæ $AC$, cū autem unaquaq; rectarū $AB$, $BC$, quæ Inter se sunt æquales sit dupla rectæ $CD$, abscindatur à rectâ $AC$, linea æqualis rectæ $CD$, et sit $CE$. Igitur recta $AE$, æquabitur $x$, et præterea tribus quartis rectæ $AC$, Igitur et $AB$, duabus tertiis ipsius $x$, et $BE$ erit tertia pars eiusdem, Igitur solidū cuius $AB$, quadratū seu $AF$, quadratū in altitudinem $BE$, erit maximum omniū quod fieri posset ex lineâ $AC$, in duas p\textsuperscript{tes} divisâ, sed eiusmodi est æquale eo quod fit sub $AB$, $BC$ et $CD$, vel $B$, $C$, $D$, hoc est $T\,\mathrm{Sol.}$ Igitur in 1\textsuperscript{o} hoc casu, $T\,\mathrm{Sol.}$ æquatur solido cuius basis quadratū duarū tertiarū ipsius $x$, altitudo verò tertia pars eiusdem \struck{\ill{}} $x$.

Sit In 2\textsuperscript{o} casu recta $CD$, minor quartâ p\textsuperscript{te} rectæ $B + C$, Igitur $AE$, vel $x$ differentia rectarū $AC$ et $CD$, erit maior tribus quartis ipsius $AC$, si Igitur $AE$ Ita dividatur In $G$, ut $AG$, sit dupla ipsius $GE$, recta $AG$ maior erit dimidiâ $AC$, et quadratū $AG$, maius rectangulo $AF$, Igitur et $GE$, etiam maior quarta p\textsuperscript{te} ipsius $AC$, ac ex consequenti maior rectâ $CD$, solidū Igitur cuius basis $AG$ quadratū, altitudo vero $GE$, maius erit solido cuius basis $AF$, quadratū altitudo vero $CD$, hoc est $T\,\mathrm{Sol.}$ cū basis sit maior basi, altitudo verò altitudine. Igitur $T\,\mathrm{Sol.}$ vel est æquale solido quod fit ex duabus tertiis $x$, in tertiam partem eiusdem, vel ipso minus.

\note{Page de droite, paginée « 17 ». Le chiffre « 20 » de la série courante dans la marge de gauche. La moitié inférieure de la page est restée blanche.}

Propositionis 5\textsuperscript{æ} Cap\textsuperscript{is} 2.\textsuperscript{i} determinatio eadem ferè quæ In superioribus In casu duorū laterū æqualiū. Cum enim In hac constitutiōe ambo latera suprà simul sumpta sint maiora reliquo quod est Infrà, et rectangulū sub summa dictorū lateriū suprà in latus infrà sit minus rectangulo sub duobus suprà, sequitur unumquodq; lateriū suprà maius esse dicto latere Infra.

Sit enim recta $AC$, summa lateriū suprà, et $AB$ sit latus $B$, et $BC$, latus $C$, recta verò $D$, sit latus infrà, rectangulū enim $AB$ in $BC$ cum sit maius ex hypothesi, rectangulo $AC$, in $D$, erit à fortiori, rectangulū $AC$ in $BC$, multo maius dicto rectangulo $AC$ in $D$ demptâ autem utrimq; $AC$, Igitur $BC$, maior erit rectâ $D$, similiq; ratiōe demonstrabitur etiam rectam $AB$, \emph{maiorem} esse ipsâ $D$.
\marginal{minorem}
\note{La figure qui accompagne ce passage est une droite lettrée $A$, $B$, $C$, avec au-dessous un segment plus court marqué $D$. Le mot « maiorem » est souligné dans le texte, et la marge de gauche porte, en regard et à côté du chiffre « 20 », le mot « minorem » : correction proposée par l'annotateur, que la démonstration de la ligne ne demande pas.}

Indè Igitur patet In hac constitutiōe unico tantū casu esse determinatiōem In quo scilicet duo latera suprà erunt Inter se æqualia, quæ quidem latera eâdem quâ in superioribus usi sumus methodo facilè elicientur.

Propositionū verò 3\textsuperscript{æ} et 4\textsuperscript{æ} Capitis determinationes eadem erunt cū superioribus, cū utraq; capita ut dictū est sint inversa superiorū, quartū scilicet primj et tertiū secundj, et in his duobus posterioribus ea dicantur de lateribus infrà, quæ in prioribus dicta sunt de lateribus suprà, si igitur laterum sic Inversorū habeatur ratio, et etiam quoad determinationes eadem dicantur de lateribus infrà, in his duobus posterioribus quæ in superioribus de lateribus Infrà, et vice versâ hic eadem de lateribus suprà quæ ibj de lateribus infrà, eadem omnino sunt determinationes.

In Capitibus autem 5\textsuperscript{o} et 6\textsuperscript{o} patet ex diversis quæ ibj sunt constitutionibus, nullas omnino esse determinationes.
\note{Une lettre « D » isolée est tracée seule sous ce dernier alinéa, au milieu de la page ; elle ressemble aux lettres de séquence qui ouvrent les parties et n'appartient pas à la phrase.}

\page{12}
\note{Double page ; page de gauche paginée « 18 », page de droite « 19 ». Le chiffre « 21. » de la série courante dans la marge de gauche.}

\subsection*{De constitutione Æquationum Quadratoquadraticarum}

Æquatio quad. quad.\textsuperscript{a} est cuius potestas est quadratoquadratū, quæ quidem æquatio cū vel a quatuor lateralibus procreetur, et hinc quinque capita, vel a quadraticâ et duabus lateralibus, vel a duabus quadraticis, vel a cubicâ (earum scilicet quæ de duobus tantū lateribus sunt explicabiles) et Laterali, vel etiam ab aliâ cubicâ, quæ per se subsistit et Laterali, vel denique hæc æquatio qq.\textsuperscript{tica} ex se subsistat, nec aliunde habeat originem, vel certè oriatur ex quibusdam lateribus, quæ veram habent significatiōem, quibusdam verò \uncertain{aliis} fictitiis simul mixtis, ideò totidem scilicet decem erunt capita, in quibus tot etiam propositiones quot varia lateriū, seu natura seu magnitudo, Inducent constitutiones.

\subsection*{Caput 1\textsuperscript{ū}}

In quo quatuor latera sunt positiva suprà.
\marginal{$S\,\mathrm{pl.pl.} - T\,\mathrm{Sol.}\,A + z\,\mathrm{pl.}\,A^{2} - xA^{3} + A^{4}$}

Propositio unica.
\[ \begin{array}{llll} \text{Sit } B - A \parallel 0 & \text{vel } C - A \parallel 0 & \text{vel } D - A \parallel 0 & \text{vel } F - A \parallel 0 \\ C - A & D - A & F - A & B - A \\ D - A & F - A & B - A & C - A \\ F - A. & B - A & C - A & D - A \end{array} \]
productū erit
\[ \begin{array}{llll} BCDF & - BCDA & + DFA^{2} & - DA^{3} \\ & - BCFA & + BFA^{2} & - BA^{3} \\ & - BDFA & + CFA^{2} & - CA^{3} \\ & - CDFA & + BCA^{2} & - FA^{3} \quad + A^{4} \quad \parallel 0 \\ & & + BDA^{2} & \\ & & + CDA^{2} & \end{array} \]
Et adhibita convenienti Interpretatione erit.
\[ S\,\mathrm{pl.pl.} - T\,\mathrm{Sol.}\,A + z\,\mathrm{pl.}\,A^{2} - xA^{3} + A^{4} \parallel 0 \]
In quâ formula sunt quatuor latera suprà sitq; $S\,\mathrm{pl.pl.}$ planoplanū sub illis; $T\,\mathrm{Sol.}$ summa triū solidorū sub iis quatuor lateribus, tribus et tribus utcumq; sumptis contentorū. $z\,\mathrm{pl.}$ aggregatū sex rectangulorū ex iisdem binis et binis, et $x$, summa eorundem laterum.

\subsection*{Caput 2.\textsuperscript{um}}

In quo sunt tria latera suprà, Alterū verò Infrà
\[ \begin{array}{llll} \text{Sit } F + A \parallel 0 & \text{vel } B - A \parallel 0 & \text{vel } C - A \parallel 0 & \text{vel } D - A \parallel 0 \\ B - A & C - A & D - A & F + A \\ C - A & D - A & F + A & B - A \\ D - A & F + A & B - A & C - A \end{array} \]
productum erit
\[ \begin{array}{llll} BCDF & + BCDA & + BFA^{2} & + DA^{3}. \\ & - BCFA & + CFA^{2} & + CA^{3} \\ & - BDFA & + DFA^{2} & + BA^{3} \\ & - CDFA & - BCA^{2} & - FA^{3} \quad - A^{4} \quad \parallel 0 \\ & & - BDA^{2} & \\ & & - CDA^{2} & \end{array} \]
Ut autem eiusmodj constitutio explicetur, cū ob diversas tum signorū $+$ et $-$ notationes, tum etiam planorū et solidorū magnitudines inæquales, varia possint oriri formulæ, videndū est primo an latus Infrà, quod hic est $F$, sit maius tribus reliquis, quæ sunt supra simul

\note{Page de droite, paginée « 19 ». Le tiers inférieur est resté blanc.}

sumptis scilicet $B$, $C$, et $D$, si enim dictū $F$, maius est his tribus, manifestum est tria rectangula quorū unū latus est Idem $F$, scilicet tria quibus præficitur signū affirmationis maiora esse tribus reliquis, quæ vocabuntur negata, cum notentur signo negatiōis; nam quia ex hypothesi, $F$, maius est $B + C + D$, rectangulū sub $F$, et summâ dictorū triū lateriū maius erit quadrato eorundem, quod quidem quadratū patet esse maius tribus rectangulis negatis, igitur et affirmata erunt iis maiora, ac proinde In hypothesi quod $F$, sit maius tribus reliquis Lateribus $z\,\mathrm{pl.}$ erit cū signo $+$, secundū quam etiam hypothesim, solidū affirmatū, erit minus tribus negatis, cū et unico eorū sit minus. Hinc Igitur

\subsection*{Proposit. 1\textsuperscript{a}}

\marginal{$S\,\mathrm{pl.pl.} - T\,\mathrm{Sol.}\,A + z\,\mathrm{pl.}\,A^{2} - xA^{3} - A^{4}$}
Cum latus Infrà maius est tribus quæ sunt suprà simul sumptis.

Cuius quidem hypothesis hæc erit formula.
\[ S\,\mathrm{pl.pl.} - T\,\mathrm{Sol.}\,A + z\,\mathrm{pl.}\,A^{2} - xA^{3} - A^{4} \parallel 0 \]
Et sunt tria latera suprà, et alterū infrà, idemq; maius tribus reliquis simul sumptis, sitq; $S\,\mathrm{pl.pl.}$ planoplanū sub dictis quatuor lateribus. $T\,\mathrm{Sol.}$ differentia solidorū sub tribus lateribus suprà In latus Infrà contentorū, et eius quod fit ex dictis tribus lateribus quæ sunt suprà, $z\,\mathrm{pl.}$ differentia rectangulorum affirmatorū et negatorū, et $x$ differentia eorumdem lateriū, et $A$ fit quodlibet ex ipsis.

\subsection*{Proposit. 2\textsuperscript{a}}

\marginal{$S\,\mathrm{pl.pl.} - T\,\mathrm{Sol.}\,A + z\,\mathrm{pl.}\,A^{2} - A^{4}$}
Cum latus Infrà æquale est summæ lateriū suprà

In qua hypothesi evanescit affectio sub cubo, et eadem rectangula afirmata maiora esse negatis eâdem viâ quâ in præcedentj propositione demonstrabitur erit Igitur
\[ S\,\mathrm{pl.pl.} - T\,\mathrm{Sol.}\,A + z\,\mathrm{pl.}\,A^{2} + A^{4} \parallel 0 \]
\note{Le dernier terme est bien écrit « $+A^{4}$ » dans la ligne, « $-A^{4}$ » dans la forme recopiée en marge ; c'est la marge qui s'accorde avec la constitution. Laissé tel quel de part et d'autre.}
Et sunt quatuor latera, tria suprà, alterū, Idemq; æquale tribus reliquis simul sumptis infrà et cæt.

Si verò latus Infrà sit minus tribus reliquis suprà simul sumptis, vel rectangula afirmata maiora, vel æqualia, vel minora erunt tribus negatis, si maiora vel æqualia, manifestū est solida sub signo negationis maiora esse eo, quod sub signo affirmationis, si verò minora, poterit etiam fieri pro ratione magnitudinis ipsius lateris, ut solida sub signo negationis vel sint maiora, vel æqualia, vel minora eo cui præficitur signū affirmatiōis, ita enim illud latus infrà poterit minui, ut sit minus quovis lateriū suprà, secundū quos varios casus, tot etiam erunt propositiones. Sit Igitur.

\subsection*{Propos. 3\textsuperscript{a}}

\marginal{$S\,\mathrm{pl.pl.} - T\,\mathrm{Sol.}\,A + z\,\mathrm{pl.}\,A^{2} + xA^{3} - A^{4}$}
Cum latus Infrà minus est summa lateriū suprà, sed rectangula affirmata maiora negatis.

Patet ex his, quæ diximus, tria solida, quæ sunt sub signo negatiōis etiam In hac constitutiōe maiora esse eo quod sub signo affirmatiōis. Erit Igitur hæc Interpretatio
\[ S\,\mathrm{pl.pl.} - T\,\mathrm{Sol.}\,A + z\,\mathrm{pl.}\,A^{2} + xA^{3} - A^{4} \parallel 0 \]
Sunt tria latera suprà, alterū Infrà idemq; minus, his quæ sunt suprà simul sumptis Et cæt.

\page{13}
\note{Double page ; page de droite paginée « 21 ». La pagination de la page de gauche est cachée par un béquet épinglé sur son angle extérieur supérieur ; le béquet relevé, elle se lit « 20 » aux vues 14 et 15. Aucun chiffre de la série courante sur l'une ni l'autre page.}

Si verò In eadē constitutiōe in quâ latus Infrà, est minus reliquis tribus quæ sunt suprà simul sumptis, rectangula verò affirmata æquentur negatis, adhuc tria solida quibus præficitur signū negatiōis erunt maiora eo quod est sub signo affirmatiōis, cū enim rectangula affirmata æquentur negatis, consequens est ut $F$, latus Infrà sit maius aliquo lateriū suprà, nam si esset minus hæc rectangula nō essent Inter se æqualia, nam $BC$, maius erit $BF$, et sic de cæteris, ponatur igitur hoc latus $F$, maius latere $B$, v.g. vel alio quovis dictorū lateriū suprà, et solido cui præficitur signū affirmatiōis statuatur basis rectangulū $CD$, altitudo verò $B$, quod ex hypothesi ipso $F$, est minus, evidens est hoc solidū, multò minus esse reliquis quæ sub signo negatiōis, et similiter procedet argumentū de quovis alio latere suprà. Sit Igitur.

\subsection*{Propos. 4\textsuperscript{a}}

\marginal{$S\,\mathrm{pl.pl.} - T\,\mathrm{Sol.}\,A + xA^{3} - A^{4}$.}
Cum latus Infrà minus est summâ lateriū suprà, et rectangula affirmata æquantur negatis

Et ob hanc æqualitatem rectangulorū \ill{} patet evanescere affectiōem sub quadrato, Cuius constitutionis hæc erit formula.
\[ S\,\mathrm{pl.pl.} - T\,\mathrm{Sol.}\,A + xA^{3} - A^{4} \parallel 0. \]
\note{Entre « rectangulorū » et « patet », un mot est écrit au-dessus de la ligne et repassé ; il reste illisible.}
Quæ quidem formula explicabitur ut In præcedentibus.

\subsection*{Propos. 5\textsuperscript{a}}

\marginal{$S\,\mathrm{pl.pl.} - T\,\mathrm{Sol.}\,A - z\,\mathrm{pl.}\,A^{2} + xA^{3} - A^{4}$}
Cum Idem latus Infrà minus est eadem summâ, et rectangula affirmata minora \uncertain{sunt} negatis, sed solidū afirmatū adhuc minus negatis

Quod quidem fieri posse demonstravimus, hæc Igitur constitutio hanc Interpretatiōem accipiet
\[ S\,\mathrm{pl.pl.} - T\,\mathrm{Sol.}\,A - z\,\mathrm{pl.}\,A^{2} + xA^{3} - A^{4} \parallel 0 \]

\subsection*{Propos. 6\textsuperscript{a}}

\marginal{$S\,\mathrm{pl.pl.} - z\,\mathrm{pl.}\,A^{2} + xA^{3} - A^{4}$.}
Iisdem positis, sed solidis negatis et affirmato Æqualibus.

Igitur in hac æquatiōe evanescit affectio sub latere et erit.
\[ S\,\mathrm{pl.pl.} - z\,\mathrm{pl.}\,A^{2} + xA^{3} - A^{4} \parallel 0 \]

\subsection*{Propos. 7\textsuperscript{a}}

\marginal{$S\,\mathrm{pl.pl.} + T\,\mathrm{Sol.}\,A - z\,\mathrm{pl.}\,A^{2} + xA^{3} - A^{4}$}
Iisdem positis adhuc ut supra sed solido affirmato maiore negatis

Erit Igitur In hac constitutiōe hæc Interpretatio.
\[ S\,\mathrm{pl.pl.} + T\,\mathrm{Sol.}\,A - z\,\mathrm{pl.}\,A^{2} + xA^{3} - A^{4} \parallel 0. \]

\subsection*{Caput 3\textsuperscript{um}}

In quo duo latera suprà, duo verò Infra.
\[ \begin{array}{ll} \text{Sit } B - A \parallel 0 & \text{vel. } D + A \parallel 0 \\ C - A & F + A. \\ D + A & B - A \\ F + A. & C - A. \end{array} \]
Productū quod inde ortus est.
\[ \begin{array}{llll} BCDF & + BCDA & + BCA^{2} & + DA^{3} \\ & + BCFA & + DFA^{2} & + FA^{3} \\ & - BDFA & - BFA^{2} & - BA^{3} \\ & - CDFA. & - CFA^{2} & - CA^{3} \quad + A^{4} \quad \parallel 0 \\ & & - CDA^{2} & \\ & & - BDA^{2} & \end{array} \]

\note{Page de droite, paginée « 21 ». Le tiers inférieur est resté blanc.}

Ut autem huiusmodj constitutio explicetur diversa est inquirenda lateriū magnitudo, prout enim vel maiora vel minora erunt, rectangula augentur, vel minuuntur, sicut etiam et solida secundū diversam rectangulorū magnitudinem, ac proinde ex variis seu lateriū, seu rectangulorū, seu etiam solidorū habitudinibus, diversæ orientur signorū notationes in formulis sub quibus variæ æquationum constitutiones concipientur.

Igitur ut horū rectangulorū diversæ magnitudines facilius dignoscantur, notandū est quatuor rectangula quibus præficitur signū negatiōis oriri ex ductu summæ lateriū quæ sunt Infrà, In summam eorū quæ sunt suprà si enim $D + F$ summa lateriū Infrà ducatur in $B + C$ summa Lateriū suprà productū erit $BD$. $BF$. $CD$. $CF$.
\[ \begin{array}{l} B + C \\ D + F \\ \hline -BD\ -BF \\ -CD\ -CF \end{array} \]
\note{Cette multiplication est posée dans la marge de gauche, en regard de la phrase.}

similiter et duo solida quæ sunt sub signo affirmatiōis oriuntur ex rectangulo quod fit à lateribus suprà, in altitudinem summæ eorū quæ sunt Infrà, sicut et duo solida quæ sunt sub signo negatiōis oriuntur ex rectangulo quod continetur sub lateribus infrà, in altitudinem summæ eorū quæ sunt suprà.

His prænotatis vel summa lateriū Infrà est maior summâ lateriū suprà, vel ipsi æqualis, vel ipsâ minor, (facilè autem iis ōibus casibus dignoscetur quo signo $x$, sit notandū ut supervacaneū sit admonere.

Et primo sit ea summa lateriū infrà maior summâ lateriū suprà, sitq; linea $AC$ ea summa lateriū infrà, et recta $DF$, summa lateriū suprà, et quidem In 1\textsuperscript{o} hoc casu quo summa lateriū infra maior est summâ lateriū suprà vel rectangula affirmata sunt maiora negatis, vel ipsis æqualia, vel minora, sed prius tanquam æqualia à nobis considerentur, ex hac enim hypothesi reliquorū casuū facilior erit determinatio.
\note{Deux droites lettrées sont tracées entre ces alinéas : la première, plus longue, porte $A$, $B$, $C$ ; la seconde, au-dessous et plus courte, $D$, $E$, $F$.}

Sunt Igitur rectangula $ABC$ et $DEF$ æqualia rectangulo $AC$ in $DF$ hoc est duo affirmata æqualia quatuor negatis, et positâ rectâ $AC$, maiore rectâ $DF$, Inde etiam sequetur ex superiori hypothesi rectangulū $ABC$, maius esse quarta p\textsuperscript{te} quadratj $DF$, scilicet maiore rectangulo quod fieri potest sub rectâ $DF$, utcumq; divisâ, si enim rectangulū $ABC$, non sit maius quartâ p\textsuperscript{te} quadratj $DF$, vel ipsi erit æquale, vel ipso minus, si æquale, duo simul scilicet rectangulū $ABC$, et quarta pars quadratj $DF$, non erunt maiora dimidio eiusdem quadratj $DF$, vel duplo quadratj $DE$, hoc est rectangulo $DF$ in $DE$, sed rectangulū $ABC$, unâ cū rectangulo $DEF$ vel quadrato $DE$, positū est æquale rectangulo sub $AC$ in $DF$, Igitur $ABC$ rectangulū In hac hypothesi maius est 4\textsuperscript{a} p\textsuperscript{te} quadratj $DF$. et ex consequenti idem rectangulū $ABC$, bis maius erit duobus rectangulis simul sumptis $ABC$ et $DEF$. sed et ex hypothesi ut Jam diximus, duo simul rectangula $ABC$ et $DEF$, æquantur rectangulo $AC$ in $DF$, Igitur rectangulū $ABC$, bis maius erit rectangulo $AC$ in $DF$, ergo multo magis idem rectangulū $ABC$, quater maius erit rectangulo $AC$ in $DF$.

\page{14}
\note{Double page. La page de gauche est blanche : c'est le verso du feuillet paginé « 21 ». Y est fixé, par une épingle qui traverse le papier, un petit béquet portant cinq lignes ; il est ici transcrit. Le béquet est rabattu sur l'angle supérieur de la page paginée « 20 » (vue 13, page de gauche), dont il masque le chiffre : celui-ci, « 20 », se lit sur le bord saillant du feuillet à cette vue et à la suivante. La page de droite est entièrement blanche et n'est pas transcrite.}

p. 21. \\
Duo verò rectangula quibus præficitur signū affirmationis nempe $BC$ et $DF$ unū fit ex ductu lateriū quæ sunt suprà nempe $B$ et $C$ et alterū ex ductu lateriū quæ sunt Infrà nempe $D$ et $F$
\note{Le béquet porte un caret sous la mention « p. 21. » : l'addition est appelée dans le texte de la page 21, après la multiplication de $B+C$ par $D+F$.}

\page{15}
\note{Double page. La page de gauche est la même page blanche qu'à la vue précédente, le premier béquet relevé : apparaît alors, sous lui et tenu par la même épingle, un second béquet de trois lignes, ici transcrit. Le chiffre « 20 » du feuillet précédent se lit toujours sur le bord saillant, à gauche. La page de droite est entièrement blanche et n'est pas transcrite.}

19\textsuperscript{a} \\
quomodō cognoscetur utrū hæc æquatio
\[ T\,\mathrm{Sol.} - z\,\mathrm{pl.}\,A + xA^{2} - A^{3} \parallel 0 \]
sit capitis quintj vel tertij.

\page{16}
\note{Double page ; page de gauche paginée « 22 », page de droite « 23 ». Les chiffres « 22 », « 23 », « 24 » de la série courante dans la marge de gauche de la page de gauche.}

erit Igitur maior ratio quadratj $DF$ ad rectangulū $AC$ in $DF$, quam eiusdem quadratj ad rectangulū $ABC$ quater sumptū. Sed ut quadratū $DF$ ad rectangulū $AC$ in $DF$, sic recta $DF$, ad rectam $AC$, sunt enim sub eâdem altitudine scilicet $DF$, Igitur maior erit ratio rectæ $DF$, ad rectam $AC$, quam quadratj $DF$, ad rectangulū $ABC$, quater sumptū, et ex consequenti eiusdem rectæ $DF$, ad eandem rectam $AC$, maior ratio, quam quater p\textsuperscript{tis} dictj quadratj $DF$, ad rectangulū $ABC$, simul sumptū. Sed rectangulū $DEF$, nō est maius quartâ p\textsuperscript{te} dictj quadratj $DF$, Igitur maior est ratio rectæ $DF$, ad rectam $AC$, quam rectangulj $DEF$, ad rectangulū $ABC$ ac ex consequenti solidū cuius basis idem rectangulū $ABC$, altitudo verò $DF$, maius erit solido cuius basis rectangulū $DEF$, altitudo verò $AC$, sed et facile demonstrabitur prius illud solidū maius esse \uncertain{septuplo} posterioris.

His Ita demonstratis constituendæ sunt ōes huius Capitis æquationū formulæ, ac 1\textsuperscript{o}. si summa lateriū Infrà, sit maior summâ lateriū suprà et rectangula afirmata sint etiam maiora negatis, Cum demonstratū \add{sit} iis existentibus Inter se æqualibus solida negata maiora esse affirmatis, igitur et multo maiora erunt. et hinc fit

\subsection*{Propositio prima}

\marginal{$S\,\mathrm{pl.pl.} - T\,\mathrm{Sol.}\,A + z\,\mathrm{pl.}\,A^{2} + xA^{3} + A^{4}$.}
Cum summa lateriū Infrà sit maior summâ lateriū suprà, et rectangula affirmata maiora negatis.

Cuius constitutionis hæc erit formula
\[ S\,\mathrm{pl.pl.} - T\,\mathrm{Sol.}\,A + z\,\mathrm{pl.}\,A^{2} + xA^{3} + A^{4} \parallel 0 \]
Et sunt quatuor latera duo suprà et duo infrà, eorumdemq; summa maior summâ horū quæ sunt suprà, $S\,\mathrm{pl.pl.}$ quod continetur sub his quatuor lateribus, $T\,\mathrm{Sol.}$ excessus quo duo solida quæ fiunt sub rectangulo lateriū Infrà in altitudinem summæ eorū quæ sunt suprà, superant duo solida quæ fiunt sub rectangulo lateriū suprà, in altitudinem summæ eorū quæ sunt Infrà, $z\,\mathrm{pl.}$ differentia rectangulorū affirmatorū et negatorū, $x$ verò differentia dictarū summarū, Et $A$, fit quodlibet ex quatuor lateribus.

\subsection*{Propos. 2\textsuperscript{a}}

\marginal{$S\,\mathrm{pl.pl.} - T\,\mathrm{Sol.}\,A + xA^{3} + A^{4}$.}
Cum eadem summa maior, et rectangula afirmata sunt æqualia negatis.

Quo quidem casu evanescet affectio sub quadrato. Et hæc erit formula
\[ S\,\mathrm{pl.pl.} - T\,\mathrm{Sol.}\,A + xA^{3} + A^{4} \parallel 0. \]
Quæ formula explicabitur ut superior et cæt. Si verò rectangula afirmata sint minora negatis, tunc vel solida negata erunt adhuc maiora affirmatis, vel ipsis æqualia, vel minora, cū enim demonstratū sit superius, quo casu hæc rectangula inter se sunt æqualia, tunc ea solida negata maiora esse affirmatis, evidens est sic ea rectangula negata posse minui, vel ut vel dicta solida negata sint adhuc maiora afirmatis, vel iis æqualia, vel denique minora, hinc tres propositiones.

\subsection*{Propos. 3\textsuperscript{a}}

\marginal{$S\,\mathrm{pl.pl.} - T\,\mathrm{Sol.}\,A - z\,\mathrm{pl.}\,A^{2} + xA^{3} + A^{4}$.}
Cū eadem summa maior, rectangula verò affirmata minora negatis, solida tamen negata maiora affirmatis

Cuius constitutionis hæc erit formula
\[ S\,\mathrm{pl.pl.} - T\,\mathrm{Sol.}\,A - z\,\mathrm{pl.}\,A^{2} + xA^{3} + A^{4} \parallel 0 \]

\note{Page de droite, paginée « 23 ». Le chiffre « 25 » de la série courante dans la marge de gauche. Le tiers inférieur est resté blanc.}

\subsection*{Propos. 4\textsuperscript{a}}

\marginal{$S\,\mathrm{pl.pl.} - z\,\mathrm{pl.}\,A^{2} + xA^{3} + A^{4}$}
Iisdem positis ut in superiore sed solidis Æqualibus.

Et hoc casu evanescit affectio sub latere.
\[ S\,\mathrm{pl.pl.} - z\,\mathrm{pl.}\,A^{2} + xA^{3} + A^{4} \parallel 0 \]

\subsection*{Propos. 5\textsuperscript{a}}

\marginal{$S\,\mathrm{pl.pl.} + T\,\mathrm{Sol.}\,A - z\,\mathrm{pl.}\,A^{2} + xA^{3} + A^{4}$.}
Iisdem positis, sed solida affirmata maiora negatis.

Tunc huius constitutionis hæc erit formula
\[ S\,\mathrm{pl.pl.} + T\,\mathrm{Sol.}\,A - z\,\mathrm{pl.}\,A^{2} + xA^{3} + A^{4} \parallel 0 \]

Si verò eadem summa lateriū Infrà sit æqualis summæ lateriū suprà manifestū est eo etiam casu rectangula quæ sunt sub signo negatiōis maiora esse iis quæ sub signo affirmatiōis, nam ea quæ sub signo negatiōis, oriuntur ex ductu summæ dictorū lateriū Infrà in summā lateriū suprà, ea verò quæ sub signo affirmatiōis producuntur ex solo ductu lateriū Inter se. Quod verò spectat ad solida, patet In hac hypothesi eam Inter se habere ratiōem quæ rectangulū ex dictis lateribus Infrà, ad rectangulū ex lateribus suprà, atq; ideò si rectangulū sub lateribus infrà maius sit rectangulo sub lateribus suprà solida negata maiora erunt affirmatis, sicut et vice versâ si rectangulū sub lateribus suprà sit maius eo sub lateribus infrà, solida affirmata erunt etiam maiora negatis, si denique dicta rectangula fiant Inter se æqualia quod aliter fieri nō potest, nisi unumquodq; lateriū Infrà sit æquale unicuiq; lateriū suprà, tunc et solida erunt etiam Inter se æqualia Sit Igitur

\subsection*{Propos. 6\textsuperscript{a}}

\marginal{$S\,\mathrm{pl.pl.} - T\,\mathrm{Sol.}\,A - z\,\mathrm{pl.}\,A^{2} + A^{4}$.}
Cum summæ lateriū sunt Inter se æquales sed rectangulū sub lateribus Infrà maius rectangulo sub lateribus suprà

Et hoc Casu evanescit affectio sub cubo.
\[ S\,\mathrm{pl.pl.} - T\,\mathrm{Sol.}\,A - z\,\mathrm{pl.}\,A^{2} + A^{4} \parallel 0. \]

\subsection*{Propos. 7\textsuperscript{a}}

\marginal{$S\,\mathrm{pl.pl.} - z\,\mathrm{pl.}\,A^{2} + A^{4}$}
Cum eadem summæ æquales et rectangulū sub lateribus Infrà æquale rectangulo sub lateribus suprà.

hoc casu evanescit non solū affectio sub cubo, sed et sub latere.
\[ S\,\mathrm{pl.pl.} - z\,\mathrm{pl.}\,A^{2} + A^{4} \parallel 0. \]

\subsection*{Propos. 8\textsuperscript{a}}

\marginal{$S\,\mathrm{pl.pl.} + T\,\mathrm{Sol.}\,A - z\,\mathrm{pl.}\,A^{2} + A^{4}$}
Cum eadem summæ lateriū sunt æquales sed rectangulū sub lateribus Infrà est minus rectangulo sub lateribus suprà.
\[ S\,\mathrm{pl.pl.} + T\,\mathrm{Sol.}\,A - z\,\mathrm{pl.}\,A^{2} + A^{4} \parallel 0 \]

Si denique summa lateriū Infrà sit minor summâ lateriū suprà, eadem dicenda erunt de lateribus suprà, quæ In præcedentibus diximus de lateribus Infrà, et similiter rectangula negata vel erant maiora, vel æqualia, vel minora afirmatis, eodemq; ratiocinio solidū cuius basis rectangulū sub lateribus suprà, altitudo verò summa lateriū infrà, maius erit solido cuius basis rectangulū sub lateribus infrà, altitudo verò, summa lateriū suprà.

\page{17}
\note{Double page ; page de gauche paginée « 24 », page de droite « 25 ».}

Et 1\textsuperscript{o}. si rectangula negata sint minora affirmatis fit.

\subsection*{Propositio 9\textsuperscript{a}}

\marginal{$S\,\mathrm{pl.pl.} + T\,\mathrm{Sol.}\,A + z\,\mathrm{pl.}\,A^{2} - xA^{3} + A^{4}$.}
Cum summa lateriū suprà maior summâ lateriū Infrà, et rectangula negata minora affirmatis.

Et tunc erit hæc formula
\[ S\,\mathrm{pl.pl.} + T\,\mathrm{Sol.}\,A + z\,\mathrm{pl.}\,A^{2} - xA^{3} + A^{4} \parallel 0. \]
Et sunt quatuor latera duo suprà et duo infrà, summa eorū quæ sunt suprà maior summâ eorū quæ sunt infrà Et cæt. ut suprà.

\subsection*{Propositio 10\textsuperscript{a}}

\marginal{$S\,\mathrm{pl.pl.} + T\,\mathrm{Sol.}\,A - xA^{3} + A^{4}$}
Iisdem positis, sed rectangulis Inter se æqualibus
\[ S\,\mathrm{pl.pl.} + T\,\mathrm{Sol.}\,A - xA^{3} + A^{4} \parallel 0 \]
Et hoc casu evanescit affectio sub quadrato.

\subsection*{Propositio 11\textsuperscript{a}}

\marginal{$S\,\mathrm{pl.pl.} + T\,\mathrm{Sol.}\,A - z\,\mathrm{pl.}\,A^{2} - xA^{3} + A^{4}$}
Cum eadem summa maior, sed rectangula affirmata minora sunt negatis, sed solida affirmata adhuc maiora.
\[ S\,\mathrm{pl.pl.} + T\,\mathrm{Sol.}\,A - z\,\mathrm{pl.}\,A^{2} - xA^{3} + A^{4} \parallel 0 \]

\subsection*{Propositio 12\textsuperscript{a}}

\marginal{$S\,\mathrm{pl.pl.} - z\,\mathrm{pl.}\,A^{2} - xA^{3} + A^{4}$.}
Iisdem positis sed solidis Æqualibus
\[ S\,\mathrm{pl.pl.} - z\,\mathrm{pl.}\,A^{2} - xA^{3} + A^{4} \parallel 0 \]
Et hoc casu evanescit affectio sub latere.

\subsection*{Propositio 13\textsuperscript{a}}

\marginal{$S\,\mathrm{pl.pl.} - T\,\mathrm{Sol.}\,A - z\,\mathrm{pl.}\,A^{2} - xA^{3} + A^{4}$}
Iisdem positis sed solidis affirmatis minoribus negatis
\[ S\,\mathrm{pl.pl.} - T\,\mathrm{Sol.}\,A - z\,\mathrm{pl.}\,A^{2} - xA^{3} + A^{4} \parallel 0 \]

\subsection*{Caput Quartum}

In quo tria latera Infrà et alterum suprà
\[ \begin{array}{llll} \text{Sit } C + A \parallel 0 & \text{vel } D + A \parallel 0 & \text{vel } F + A \parallel 0 & \text{vel } B - A \\ D + A & F + A & B - A & C + A \\ F + A & B - A & C + A & D + A \\ B - A & C + A & D + A & F + A \end{array} \]
Productum autem erit
\[ \begin{array}{llll} BCDF & + BDFA & + BCA^{2} & + BA^{3} \\ & + BCFA & + BDA^{2} & - CA^{3} \\ & + BCDA & + BFA^{2} & - DA^{3} \\ & - DCFA & - CDA^{2} & - FA^{3} \quad - A^{4} \quad \parallel 0 \\ & & - CFA^{2}. & \\ & & - DFA^{2} & \end{array} \]
Quæ constitutio ut explicetur notanda et observanda sunt ea quæ in 2\textsuperscript{o} huius p\textsuperscript{tis} capite adnotata sunt, modo hic Intelligantur de lateribus infrà, quæ ibj dicta sunt de lateribus suprà, similiter hic de lateribus suprà quæ illic de lateribus Infrà, unde et totidem erunt huius Capitis propositiones.

\note{Page de droite, paginée « 25 ». La moitié inférieure est restée blanche.}

\subsection*{Propositio 1\textsuperscript{a}}

\marginal{$S\,\mathrm{pl.pl.} + T\,\mathrm{Sol.}\,A + z\,\mathrm{pl.}\,A^{2} + xA^{3} - A^{4}$}
Cum latus supra maius est tribus reliquis simul sumptis
\[ S\,\mathrm{pl.pl.} + T\,\mathrm{Sol.}\,A + z\,\mathrm{pl.}\,A^{2} + xA^{3} - A^{4} \parallel 0 \]
Et sunt quatuor latera tria infrà, et alterū Idemq; maius reliquis simul sumptis suprà, $S\,\mathrm{pl.pl.}$ quod fit sub his lateribus, $T\,\mathrm{Sol.}$ differentia solidorū et cæt. ut In dicto 2\textsuperscript{o} cap.

\subsection*{Propositio 2\textsuperscript{a}}

\marginal{$S\,\mathrm{pl.pl.} + T\,\mathrm{Sol.}\,A + z\,\mathrm{pl.}\,A^{2} - A^{4}$.}
Cum Idem latus suprà æquale est iis quæ sunt supra.
\[ S\,\mathrm{pl.pl.} + T\,\mathrm{Sol.}\,A + z\,\mathrm{pl.}\,A^{2} - A^{4} \parallel 0 \]
Et evanescit In hoc casu affectio sub latere.
\note{« quæ sunt supra » est bien ce qui est écrit ; le chapitre pose trois côtés \emph{infra} et un seul \emph{supra}, et c'est « infra » que la proposition demande. Laissé tel quel.}

\subsection*{Propositio 3\textsuperscript{a}}

\marginal{$S\,\mathrm{pl.pl.} + T\,\mathrm{Sol.}\,A + z\,\mathrm{pl.}\,A^{2} - xA^{3} - A^{4}$}
Cum Idem latus supra minus est tribus iis, quæ sunt Infrà simul sumptis, sed rectangula affirmata maiora negatis
\[ S\,\mathrm{pl.pl.} + T\,\mathrm{Sol.}\,A + z\,\mathrm{pl.}\,A^{2} - xA^{3} - A^{4} \parallel 0 \]

\subsection*{Propositio 4\textsuperscript{a}}

\marginal{$S\,\mathrm{pl.pl.} + T\,\mathrm{Sol.}\,A - xA^{3} - A^{4}$}
Iisdem positis sed rectangulis æqualibus
\[ S\,\mathrm{pl.pl.} + T\,\mathrm{Sol.}\,A - xA^{3} - A^{4} \parallel 0 \]
Et hoc casu evanescit affectio sub quadrato.

\subsection*{Propositio 5\textsuperscript{a}}

\marginal{$S\,\mathrm{pl.pl.} + T\,\mathrm{Sol.}\,A - z\,\mathrm{pl.}\,A^{2} - xA^{3} - A^{4}$}
Cum Idem latus Infrà minus est, rectangula negata etiam minora afirmatis, sed solida affirmata adhuc maiora negato.
\[ S\,\mathrm{pl.pl.} + T\,\mathrm{Sol.}\,A - z\,\mathrm{pl.}\,A^{2} - xA^{3} - A^{4} \parallel 0 \]

\subsection*{Propositio 6\textsuperscript{a}}

\marginal{$S\,\mathrm{pl.pl.} - z\,\mathrm{pl.}\,A^{2} - xA^{3} - A^{4}$.}
Iisdem positis sed solidis æqualibus
\[ S\,\mathrm{pl.pl.} - z\,\mathrm{pl.}\,A^{2} - xA^{3} - A^{4} \parallel 0 \]
Et hoc casu evanescit affectio sub latere.

\subsection*{Propositio 7\textsuperscript{a}}

\marginal{$S\,\mathrm{pl.pl.} - T\,\mathrm{Sol.}\,A - z\,\mathrm{pl.}\,A^{2} - xA^{3} - A^{4}$.}
Cum Idem latus est minus rectangula affirmata minora negatis, et solida etiam affirmata minora negatis
\[ S\,\mathrm{pl.pl.} - T\,\mathrm{Sol.}\,A - z\,\mathrm{pl.}\,A^{2} - xA^{3} - A^{4} \parallel 0 \]

\subsection*{Caput Quintum}

In quo omnia latera sunt Infrà.
\[ \begin{array}{llll} \text{Sit } B + A \parallel 0 & \text{vel } C + A \parallel 0 & \text{vel } D + A \parallel 0 & \text{vel } F + A \parallel 0 \\ C + A & D + A & F + A & B + A \\ D + A & F + A & B + A & C + A \\ F + A & B + A. & C + A & D + A \end{array} \]
\marginal{$S\,\mathrm{pl.pl.} + T\,\mathrm{Sol.}\,A + z\,\mathrm{pl.}\,A^{2} + xA^{3} + A^{4}$.}
Et adhibita convenienti Interpretatione erit
\[ S\,\mathrm{pl.pl.} + T\,\mathrm{Sol.}\,A + z\,\mathrm{pl.}\,A^{2} + xA^{3} + A^{4} \parallel 0 \]
Et cū hoc caput sit primj Inversum, sic et formula erit omnino inversa, et ut explicetur hic dicendū est de lateribus infrà quæ ibj dicta sunt de iis quæ sunt supra.

\page{18}
\note{Double page ; page de gauche paginée « 26 », page de droite « 27 ». Le chiffre « 26 » de la série courante dans la marge de gauche de la page de gauche.}

\subsection*{Caput 6\textsuperscript{um}. De Constitutione Æquationū quadratoquadraticarum quæ oriuntur ex duabus planis}

Cum In eâ parte, in quâ superius egimus de æquationibus quadratis dictū sit eas esse tantummodo duplicis generis, vel quæ ex duobus lateribus, vel quæ ex unico tantū explicantur, hic nō est quæstio de æquationibus qq\textsuperscript{tis} quæ originem ex duabus quadratis, quarū unaquæq; duas radices, eo enim casu productū quod inde procreatur sub aliquâ ex superioribus huius partis formulis, quæ ex quatuor lateribus explicantur, interpretandū esset, hic Igitur agitur de iis quadratis quarū unicū est latus, cū verò ex cognitiōe ipsarū manifestū sit, esse tantū unicam huiusmodi æquationū, Igitur huius capitis erit

\subsection*{Propositio unica}

In qua duo latera ambo Infra.
\[ \begin{array}{l} \text{Sit enim } G\,\mathrm{pl.} + A^{2} \parallel 0. \\ \text{et } H\,\mathrm{pl.} + A^{2} \parallel 0 \end{array} \]
productū quod Inde oritur erit.
\[ G\,\mathrm{pl.}\,H\,\mathrm{pl.} + G\,\mathrm{pl.}\,A^{2} + H\,\mathrm{pl.}\,A^{2} + A^{4} \parallel 0 \]
\note{Les deux termes en $A^{2}$ sont écrits l'un au-dessus de l'autre après « $G\,\mathrm{pl.}\,H\,\mathrm{pl.}$ » ; ils sont ici mis bout à bout.}
Quæ constitutio hanc accipit Interpretationem.
\marginal{$S\,\mathrm{pl.pl.} + z\,\mathrm{pl.}\,A^{2} + A^{4}$.}
\[ S\,\mathrm{pl.pl.} + z\,\mathrm{pl.}\,A^{2} + A^{4} \parallel 0 \]
Et sunt duo latera Infrà quorū quadrata simul sumpta exhibent $z\,\mathrm{pl.}$ et $S\,\mathrm{pl.pl.}$ quod sub his quadratis continetur. $A$, verò fit utrumlibet ex his duobus lateribus.

\subsection*{Caput 7\textsuperscript{um}. De constitutione æquationū quadratoquadraticarū quæ oriuntur ex æquatione planâ et duabus lateralibus}

Ex superiori capite patet de quâ æquatiōe quadratâ hic agatur, duæ verò laterales vel enunciabuntur de duobus lateribus suprà, vel de uno latere suprà, altero Infrà, vel de duobus Infrà.
\[ \text{Sint } 1^{\mathrm{o}} \text{ duo latera suprà} \quad \begin{array}{l} B - A \parallel 0. \\ C - A \parallel 0. \end{array} \]
productū quod Inde oritur, si ducatur in dictam æquatiōem quadratam erit
\[ \begin{array}{llll} BCG\,\mathrm{pl.} & - BG\,\mathrm{pl.}\,A & + G\,\mathrm{pl.}\,A^{2} & - BA^{3} \\ & - CG\,\mathrm{pl.}\,A & + BCA^{2} & - CA^{3} \quad + A^{4} \quad \parallel 0 \end{array} \]
Hinc Igitur

\subsection*{Propositio 1\textsuperscript{a}}

\marginal{$S\,\mathrm{pl.pl.} - T\,\mathrm{Sol.}\,A + z\,\mathrm{pl.}\,A^{2} - xA^{3} + A^{4}$.}
Cum ex planâ et duabus lateralibus quarū lateraliū radices sunt suprà

Cui constitutionj hæc erit Interpretatio.
\[ S\,\mathrm{pl.pl.} - T\,\mathrm{Sol.}\,A + z\,\mathrm{pl.}\,A^{2} - xA^{3} + A^{4} \parallel 0 \]
Quæ formula etiam si similis sit ej quæ habetur in Cap\textsuperscript{ite} 1\textsuperscript{o} huius p\textsuperscript{tis}, maximè tamen Inter se differunt, ut videre est ex earū diversâ originè nam hic sunt tantum tria latera, duo suprà, alterū verò Infra

\note{Page de droite, paginée « 27 ». Le tiers inférieur est resté blanc.}

$x$ est summa eorū quæ sunt suprà, summa rectangulj quod sub iisdem continetur et quadratj eius quod est Infrà est $z\,\mathrm{pl.}$, et $T\,\mathrm{Sol.}$ fit ex dicto quadrato in altitudinem summæ dictorū lateriū quæ sunt suprà, Et $S\,\mathrm{pl.pl.}$ quod sub eodem quadrato et dicto rectangulo continetur, $A$ verò fit utrumlibet dictorū lateriū

\subsection*{Propositio 2\textsuperscript{a}}

Cum iisdem positis, una ex radicibus lateraliū est suprà, altera verò infrà, et ea quæ est suprà maior est eâ quæ est Infrà, et rectangulū sub dictis radicibus maius plano æquationis quadratæ.

Huiusmodj constitutio inde oritur.
\[ \begin{array}{ll} \text{Sit } B - A \parallel 0 & \text{productū erit} \\ \text{et } C + A \parallel 0 & BC + BA - CA - A^{2} \parallel 0 \end{array} \]
Quod productū si ducatur in $G^{2} + A^{2} \parallel 0$ erit
\[ \begin{array}{llll} BCG\,\mathrm{pl.} & + BG\,\mathrm{pl.}\,A & - G\,\mathrm{pl.}\,A^{2} & + BA^{3} \\ & - CG\,\mathrm{pl.}\,A & + BCA^{2} & - CA^{3} \quad - A^{4} \quad \parallel 0 \end{array} \]
\note{Le plan de l'équation quadratique, écrit « $G\,\mathrm{pl.}$ » partout ailleurs, est ici noté « $G^{2}$ » ; conservé tel quel.}

Quoniam verò fieri potest ut latus suprà sit vel maius, vel minus latere infrà, vel ipsi æquale, et similiter rectangulū $BC$, vel maius, vel æquale, vel minus $G\,\mathrm{pl.}$, quod et Idem etiam potest accidere circa solida, Ideo in casu huiusce 1\textsuperscript{æ} propositionis hæc sit formula
\marginal{$S\,\mathrm{pl.pl.} + T\,\mathrm{Sol.}\,A + z\,\mathrm{pl.}\,A^{2} + xA^{3} - A^{4}$.}
\[ S\,\mathrm{pl.pl.} + T\,\mathrm{Sol.}\,A + z\,\mathrm{pl.}\,A^{2} + xA^{3} - A^{4} \parallel 0 \]
Nec dissimilis est Istius formulæ explicatio à, præcedentj, nisi quod \struck{et} quæ in superiorj erant summæ hic sunt differentiæ.

\subsection*{Propositio 3.\textsuperscript{a}}

\marginal{$S\,\mathrm{pl.pl.} + T\,\mathrm{Sol.}\,A + xA^{3} - A^{4}$}
Cum Idem latus suprà est maius, sed rectangulū sub lateribus æquale $G\,\mathrm{pl.}$
\[ S\,\mathrm{pl.pl.} + T\,\mathrm{Sol.}\,A + xA^{3} - A^{4} \parallel 0 \]
Et hoc casu evanescit affectio sub quadrato.

\subsection*{Propositio 4\textsuperscript{a}}

\marginal{$S\,\mathrm{pl.pl.} + T\,\mathrm{Sol.}\,A - z\,\mathrm{pl.}\,A^{2} + xA^{3} - A^{4}$}
Cum latus supra est maius et rectangulū sub lateribus est \add{minus} $G\,\mathrm{pl.}$
\[ S\,\mathrm{pl.pl.} + T\,\mathrm{Sol.}\,A - z\,\mathrm{pl.}\,A^{2} + xA^{3} - A^{4} \parallel 0 \]
\note{« minus » est ajouté au-dessus de la ligne, appelé par un caret.}

\subsection*{Propositio 5\textsuperscript{a}}

\marginal{$S\,\mathrm{pl.pl.} + z\,\mathrm{pl.}\,A^{2} - A^{4}$}
Cum latera sunt æqualia, sed rectangulū sub iis lateribus est maius $G\,\mathrm{pl.}$
\[ S\,\mathrm{pl.pl.} + z\,\mathrm{pl.}\,A^{2} - A^{4} \parallel 0 \]
Et hoc casu evanescit affectio sub latere et sub cubo, utriusq; enim solidj eadem est basis.

\subsection*{Propositio 6\textsuperscript{a}}

\marginal{$S\,\mathrm{pl.pl.} - A^{4}$}
Iisdem positis, et rectangulo æqualj $G\,\mathrm{pl.}$
\[ S\,\mathrm{pl.pl.} - A^{4} \parallel 0 \]
Et omnes affectiones sub gradibus evanescunt.

\page{19}
\note{Double page ; page de gauche paginée « 28. », page de droite « 29. ». Le chiffre « 27. » de la série courante dans la marge de gauche de la page de gauche, « 29 » dans celle de la page de droite. Au-dessus de l'angle extérieur supérieur de la page de droite dépasse le bord d'un feuillet postérieur, qui porte en gros chiffres « 33 » ; c'est sa pagination, et non un chiffre de celui-ci.}

\subsection*{Propositio 7\textsuperscript{a}}

\marginal{$S\,\mathrm{pl.pl.} - z\,\mathrm{pl.}\,A^{2} - A^{4}$.}
Cum eadem latera æqualia, sed rectangulo sub iis minore $G\,\mathrm{pl.}$
\[ S\,\mathrm{pl.pl.} - z\,\mathrm{pl.}\,A^{2} - A^{4} \parallel 0 \]

\subsection*{Propositio 8\textsuperscript{a}}

\marginal{$S\,\mathrm{pl.pl.} - T\,\mathrm{Sol.}\,A + z\,\mathrm{pl.}\,A^{2} - xA^{3} - A^{4}$.}
Cum latus Infrà est maius latere suprà et Rectangulo sub lateribus maiori $G\,\mathrm{pl.}$
\[ S\,\mathrm{pl.pl.} - T\,\mathrm{Sol.}\,A + z\,\mathrm{pl.}\,A^{2} - xA^{3} - A^{4} \parallel 0 \]

\subsection*{Propositio 9\textsuperscript{a}}

\marginal{$S\,\mathrm{pl.pl.} - T\,\mathrm{Sol.}\,A - xA^{3} - A^{4}$.}
Iisdem positis et rectangulo æqualj $G\,\mathrm{pl.}$
\[ S\,\mathrm{pl.pl.} - T\,\mathrm{Sol.}\,A - xA^{3} - A^{4} \parallel 0 \]
Et evanescit affectio sub quadrato.

\subsection*{Propositio 10\textsuperscript{a}}

\marginal{$S\,\mathrm{pl.pl.} - T\,\mathrm{Sol.}\,A - z\,\mathrm{pl.}\,A^{2} - xA^{3} - A^{4}$}
Cum Idem latus Infra est maius et rectangulō sub lateribus minore $G\,\mathrm{pl.}$
\[ S\,\mathrm{pl.pl.} - T\,\mathrm{Sol.}\,A - z\,\mathrm{pl.}\,A^{2} - xA^{3} - A^{4} \parallel 0 \]

Si denique hæc duo latera sint ambo Infra
\[ \begin{array}{ll} \text{Sit } B + A \parallel 0 & \text{productū erit} \\ \text{et } C + A \parallel 0. & BC + BA + CA + A^{2} \parallel 0 \end{array} \]
Quod productū si ducatur ut suprà in $G\,\mathrm{pl.} + A^{2} \parallel 0$ Erit et Istud
\[ \begin{array}{llll} BCG\,\mathrm{pl.} & + BG\,\mathrm{pl.}\,A & + G\,\mathrm{pl.}\,A^{2} & + BA^{3} \\ & + CG\,\mathrm{pl.}\,A & + BCA^{2} & + CA^{3} \quad + A^{4} \quad \parallel 0 \end{array} \]
Ex quo

\subsection*{Propositio 11\textsuperscript{a}}

\marginal{$S\,\mathrm{pl.pl.} + T\,\mathrm{Sol.}\,A + z\,\mathrm{pl.}\,A^{2} + xA^{3} + A^{4}$}
Cum ambo latera sunt Infra.
\[ S\,\mathrm{pl.pl.} + T\,\mathrm{Sol.}\,A + z\,\mathrm{pl.}\,A^{2} + xA^{3} + A^{4} \parallel 0 \]
Et sunt tria latera Infrà et cæt. ut In 1\textsuperscript{a} huius Capitis modo diversorū lateriū habeatur ratio.

\subsection*{Caput 8.\textsuperscript{um} De constitutione Æquationū quad. quad.\textsuperscript{arū} quæ oriuntur ex Æquatione cubica et lateralj}

Eadem In hoc capite et In sequentj erunt adnotanda quæ superius In capite \uncertain{sexto} huius p\textsuperscript{tis}, scilicet hic nō agi de iis æquationibus cubicis quæ de tribus lateribus explicantur, sed tantū de iis quæ vel ex duobus, vel etiam de iis quæ per se subsistunt, quod in Cap\textsuperscript{te} præfato, si enim illæ æquationes cubicæ quarū duæ sunt radices ducantur in lateralem oriuntur quad. quadrata de tribus lateribus tantummodo explicabiles.

Harū verò æquationū Cubicarū cū sint tantū duæ formulæ, ut patet ex recognitiōe ipsarū, una In quâ alterū ex lateribus est suprà alterū verò infrà, alia verò in quâ ambo latera sunt infrà, quarū utraq; potest ducj, vel in latus suprà, vel in latus Infrà. Si Igitur prima formula quam sic Interpretatj sumus.

\note{Page de droite, paginée « 29. ». Le chiffre « 29 » de la série courante dans la marge de gauche. La moitié inférieure est restée blanche.}

\[ P\,\mathrm{Sol.} - G\,\mathrm{pl.}\,A + FA^{2} - A^{3} \parallel 0 \]
ducatur primo in lateralem Cuius latus suprà scilicet $B - A \parallel 0$ productum erit
\[ \begin{array}{llll} BP\,\mathrm{Sol.} & - BG\,\mathrm{pl.}\,A & + G\,\mathrm{pl.}\,A^{2} & - BA^{3} \\ & - P\,\mathrm{Sol.}\,A & + BFA^{2} & - FA^{3} \quad - A^{4} \quad \parallel 0 \end{array} \]
Et adhibita convenienti Interpretatiōe, erit

\subsection*{Propositio 1\textsuperscript{a}}

\marginal{$S\,\mathrm{pl.pl.} - T\,\mathrm{Sol.}\,A + z\,\mathrm{pl.}\,A^{2} - xA^{3} - A^{4}$}
Cum duo latera sunt suprà, alterū Infra.
\[ S\,\mathrm{pl.pl.} - T\,\mathrm{Sol.}\,A + z\,\mathrm{pl.}\,A^{2} - xA^{3} - A^{4} \parallel 0 \]
Sunt Igitur tria latera duo suprà, alterū Infrà, $x$, est summa eorū quæ sunt suprà, $z\,\mathrm{pl.}$ aggregatū rectangulj quod continetur sub illis, et quadratj eius quod est infrà, $T\,\mathrm{Sol.}$, solidū quod fit ex quadrato dictj lateris infra in altitudinem summæ eorū quæ sunt suprà. Et $S\,\mathrm{pl.pl.}$ quod sub dictis quadrato et rectangulo.

Si verò eadem formula ducatur in lateralem $C + A \parallel$ Cuius latus Infra productū erit
\[ \begin{array}{llll} P\,\mathrm{Sol.}\,C & + P\,\mathrm{Sol.}\,A & + CFA^{2} & + FA^{3} \\ & - CG\,\mathrm{pl.}\,A & - G\,\mathrm{pl.}\,A^{2} & - CA^{3} \quad - A^{4} \quad \parallel 0 \end{array} \]
Ex quo patet ob varias lateriū magnitudines, et signorū diversas notationes varias etiam oriri formulas, vel enim latus infra $C$, maius est latere suprà $F$, vel ipsi æquale, vel ipso minus, sic etiam eodem vel maiore vel æquali, vel minore dicto latere suprà fieri potest ut circa rectangula eadem accidant, quod enim spectat ad solida sequuntur variam horū lateriū naturam et magnitudinem. Hinc Igitur

\subsection*{Propositio 2\textsuperscript{a}}

\marginal{$S\,\mathrm{pl.pl.} - T\,\mathrm{Sol.}\,A + z\,\mathrm{pl.}\,A^{2} - xA^{3} + A^{4}$}
Cum latus Infrà maius est eo quod est suprà, sed rectangulo sub his lateribus maiore quadrato eius quod est potentiâ æquale $G\,\mathrm{Pl.}$
\[ S\,\mathrm{pl.pl.} - T\,\mathrm{Sol.}\,A + z\,\mathrm{pl.}\,A^{2} - xA^{3} + A^{4} \parallel 0 \]
Suntq; tria latera duo infrà (quorū unū æquatur potentiâ $G\,\mathrm{pl.}$ alterū verò est $C$) et aliud supra quod est $F$, latus scilicet æquationis cubicæ, idemq; minus dicto latere $C$, quod est Infrà, $x$, est differentia dictorū lateriū $C$ et $F$. Et $z\,\mathrm{pl.}$ differentia rectangulj sub iisdem, et quadratj alterius æquationis cubicæ quod est Infrà, $T\,\mathrm{Sol.}$ differentia solidorū quorū basis Idem quadratū, altitudo verò, unius scilicet $C$, alterius verò $F$, et $S\,\mathrm{pl.pl.}$ quod continetur sub his rectangulo et quadrato.

\subsection*{Propositio 3\textsuperscript{a}}

\marginal{$S\,\mathrm{pl.pl.} - T\,\mathrm{Sol.}\,A - xA^{3} + A^{4}$}
Cum Idem latus Infrà est maius, sed rectangulū sub dictis lateribus æquale quadrato alterius.
\[ S\,\mathrm{pl.pl.} - T\,\mathrm{Sol.}\,A - xA^{3} + A^{4} \parallel 0 \]
Cuius formulæ \struck{et} sicut et sequentiū facilis est explicatio ex superiorj.

\page{20}
\note{Double page ; page de gauche paginée « 30 », page de droite « 31. ». Le bord du feuillet postérieur paginé « 33 » dépasse encore au-dessus de l'angle extérieur de la page de droite. Aucun chiffre de la série courante.}

\subsection*{Propositio 4\textsuperscript{a}}

\marginal{$S\,\mathrm{pl.pl.} - T\,\mathrm{Sol.}\,A - z\,\mathrm{pl.}\,A^{2} - xA^{3} + A^{4}$}
Cum Idem latus est maius, sed dictū rectangulū minus eodem quad\textsuperscript{to}.
\[ S\,\mathrm{pl.pl.} - T\,\mathrm{Sol.}\,A - z\,\mathrm{pl.}\,A^{2} - xA^{3} + A^{4} \parallel 0 \]

\subsection*{Propositio 5\textsuperscript{a}}

\marginal{$S\,\mathrm{pl.pl.} - z\,\mathrm{pl.}\,A^{2} + A^{4}$.}
Cum eadem latera Inter se æqualia, sed rectangulū sub iis maius dicto quadrato.
\[ S\,\mathrm{pl.pl.} + z\,\mathrm{pl.}\,A^{2} + A^{4} \parallel 0 \]
\note{La ligne porte « $+\,z\,\mathrm{pl.}\,A^{2}$ », la forme recopiée en marge « $-\,z\,\mathrm{pl.}\,A^{2}$ » ; les deux sont donnés tels quels.}

\subsection*{Propositio 6\textsuperscript{a}}

\marginal{$S\,\mathrm{pl.pl.} + A^{4}$.}
Iisdem positis, et rectangulo æqualj quadrato.
\[ S\,\mathrm{pl.pl.} + A^{4} \parallel 0 \]

\subsection*{Propositio 7\textsuperscript{a}}

\marginal{$S\,\mathrm{pl.pl.} - z\,\mathrm{pl.}\,A^{2} + A^{4}$.}
Iisdem lateribus Inter se Æqualibus, sed Idem rectangulum minus dicto quadrato.
\[ S\,\mathrm{pl.pl.} - z\,\mathrm{pl.}\,A^{2} + A^{4} \parallel 0 \]

\subsection*{Propositio 8\textsuperscript{a}}

\marginal{$S\,\mathrm{pl.pl.} + T\,\mathrm{Sol.}\,A + z\,\mathrm{pl.}\,A^{2} + xA^{3} + A^{4}$}
Cum Idem latus Infrà minus est dicto latere suprà, sed dictū Rectangulū maius dicto quadrato
\[ S\,\mathrm{pl.pl.} + T\,\mathrm{Sol.}\,A + z\,\mathrm{pl.}\,A^{2} + xA^{3} + A^{4} \parallel 0 \]

\subsection*{Propositio 9\textsuperscript{a}}

\marginal{$S\,\mathrm{pl.pl.} + T\,\mathrm{Sol.}\,A + xA^{3} + A^{4}$}
Eodem latere minore, sed rectangulo et quadrato æqualibus
\[ S\,\mathrm{pl.pl.} + T\,\mathrm{Sol.}\,A + xA^{3} + A^{4} \parallel 0 \]

\subsection*{Propositio 10\textsuperscript{a}}

\marginal{$S\,\mathrm{pl.pl.} + T\,\mathrm{Sol.}\,A - z\,\mathrm{pl.}\,A^{2} + xA^{3} + A^{4}$}
Eodem latere minore et dicto rectangulo minore dicto quadrato.
\[ S\,\mathrm{pl.pl.} + T\,\mathrm{Sol.}\,A - z\,\mathrm{pl.}\,A^{2} + xA^{3} + A^{4} \parallel 0 \]

Si deinde formula alterius æquationis Cubicæ scilicet
\[ P.\,\mathrm{Sol} + G\,\mathrm{pl.}\,A + FA^{2} + A^{3} \parallel 0 \]
In qua sunt duo latera Infrà ducatur in lateralem, radix huius lateralis erit ut in præcedentj Constitutiōe vel suprà vel Infrà, sit 1\textsuperscript{o} suprà scilicet $B - A \parallel 0$ productū erit huiusmodj.
\[ \begin{array}{llll} BP\,\mathrm{pl.} & + BG\,\mathrm{pl.}\,A & + BFA^{2} & + BA^{3} \\ & - P\,\mathrm{Sol.}\,A & - G\,\mathrm{pl.}\,A^{2} & - FA^{3} \quad - A^{4} \quad \parallel 0 \end{array} \]
\note{Le produit est précédé d'un mot biffé, illisible sous la rature, et écrit en deux lignes accolées ; elles sont ici mises bout à bout.}
Ex quo producto patet ut ex superiorj, varias oriri formulas, ob diversas signorū notationes, idq; pro variâ lateriū magnitudine, quæ quidem formulæ etsi similes erunt fere ōibus superioribus huius Capitis, earū diversa est explicatio, ut patet ex ipsarū origine.

\subsection*{Propositio 11\textsuperscript{a}}

\marginal{$S\,\mathrm{pl.pl.} + T\,\mathrm{Sol.}\,A + z\,\mathrm{pl.}\,A^{2} + xA^{3} - A^{4}$}
Cū latus suprà maius est latere infrà æquationis cubicæ, et rectangulū sub his lateribus maius, quadrato alterius lateris Infrà eiusdem æquationis Cubicæ
\[ S\,\mathrm{pl.pl.} + T\,\mathrm{Sol.}\,A + z\,\mathrm{pl.}\,A^{2} + xA^{3} - A^{4} \parallel 0 \]
Sunt tria latera, duo Infrà, quæ spectant ad æquationem cubicam, alterū quidem æquale potentiâ $G\,\mathrm{pl.}$ reliquū verò eiusdem æquationis est $F$, tertiū

\note{Page de droite, paginée « 31. ». Le tiers inférieur est resté blanc. C'est la dernière page de cette livraison ; le texte se poursuit à la vue 21.}

verò latus in hac formulâ est suprà, idemq; maius dicto latere $F$, differentia horū lateriū est $x$, differentia rectangulj sub his et dictj quadratj $z\,\mathrm{pl.}$ et cæt. ut supra

\subsection*{Propositio 12\textsuperscript{a}}

\marginal{$S\,\mathrm{pl.pl.} + T\,\mathrm{Sol.}\,A + xA^{3} - A^{4}$}
Cum Idem latus supra est maius, sed Idem rectangulū æquale dicto quadrato.
\[ S\,\mathrm{pl.pl.} + T\,\mathrm{Sol.}\,A + xA^{3} - A^{4} \parallel 0 \]

\subsection*{Propositio 13\textsuperscript{a}}

\marginal{$S\,\mathrm{pl.pl.} + T\,\mathrm{Sol.}\,A - z\,\mathrm{pl.}\,A^{2} + xA^{3} - A^{4}$}
Eodem latere maiore, et dicto rectangulo minore dicto quadrato
\[ S\,\mathrm{pl.pl.} + T\,\mathrm{Sol.}\,A - z\,\mathrm{pl.}\,A^{2} + xA^{3} - A^{4} \parallel 0 \]

\subsection*{Propositio 14\textsuperscript{a}}

\marginal{$S\,\mathrm{pl.pl.} + z\,\mathrm{pl.}\,A^{2} - A^{4}$}
Iisdem lateribus $B$ et $F$ æqualibus et dicto rectangulo maiore dicto quadrato
\[ S\,\mathrm{pl.pl.} + z\,\mathrm{pl.}\,A^{2} - A^{4} \parallel 0 \]

\subsection*{Propositio 15\textsuperscript{a}}

\marginal{$S\,\mathrm{pl.pl.} - A^{4}$}
Iisdem lateribus, et dictis rectangulo et quadrato æqualibus
\[ S\,\mathrm{pl.pl.} - A^{4} \parallel 0 \]

\subsection*{Propositio 16\textsuperscript{a}}

\marginal{$S\,\mathrm{pl.pl.} - z\,\mathrm{pl.}\,A^{2} - A^{4}$}
Iisdem lateribus æqualibus et dicto rectangulo minore dicto quadrato.
\[ S\,\mathrm{pl.pl.} - z\,\mathrm{pl.}\,A^{2} - A^{4} \parallel 0 \]

\subsection*{Propositio 17\textsuperscript{a}}

\marginal{$S\,\mathrm{pl.pl.} - T\,\mathrm{Sol.}\,A + z\,\mathrm{pl.}\,A^{2} - xA^{3} - A^{4}$.}
Cū Idem latus suprà est minus dicto latere Infrà, et dictum rectangulū maius dicto quadrato
\[ S\,\mathrm{pl.pl.} - T\,\mathrm{Sol.}\,A + z\,\mathrm{pl.}\,A^{2} - xA^{3} - A^{4} \parallel 0 \]

\subsection*{Propositio 18\textsuperscript{a}}

\marginal{$S\,\mathrm{pl.pl.} - T\,\mathrm{Sol.}\,A - xA^{3} - A^{4}$}
Eodem latere minore, et dictis rectangulo et quadrato æqualibus
\[ S\,\mathrm{pl.pl.} - T\,\mathrm{Sol.}\,A - xA^{3} - A^{4} \parallel 0 \]

\subsection*{Propositio 19\textsuperscript{a}}

\marginal{$S\,\mathrm{pl.pl.} - T\,\mathrm{Sol.}\,A - z\,\mathrm{pl.}\,A^{2} - xA^{3} - A^{4}$}
Eodem latere minore et dicto rectangulo etiam minore dicto quadrato
\[ S\,\mathrm{pl.pl.} - T\,\mathrm{Sol.}\,A - z\,\mathrm{pl.}\,A^{2} - xA^{3} - A^{4} \parallel 0 \]

Si denique hæc ultima formula æquationis cubicæ ducatur In lateralem, cuius latus sit positivū Infra scilicet $C + A \parallel 0$ productū erit
\[ \begin{array}{llll} P\,\mathrm{Sol.}\,C & + CG\,\mathrm{pl.}\,A & + CFA^{2} & + CA^{3} \\ & + P\,\mathrm{Sol.}\,A & + G\,\mathrm{pl.}\,A^{2} & + FA^{3} \quad + A^{4} \quad \parallel 0 \end{array} \]
Et adhibitâ convenienti Interpretatiōe fiet.

\subsection*{Propositio 20\textsuperscript{a}}

\marginal{$S\,\mathrm{pl.pl.} + T\,\mathrm{Sol.}\,A + z\,\mathrm{pl.}\,A^{2} + xA^{3} + A^{4}$.}
Cum tria latera sunt Infra.
\[ S\,\mathrm{pl.pl.} + T\,\mathrm{Sol.}\,A + z\,\mathrm{pl.}\,A^{2} + xA^{3} + A^{4} \parallel 0 \]
Et sunt tria latera infrà, $x$ est aggregatū duorū quorū unū spectat ad dictam æquationem cubicam, alterū verò ad lateralem propositam, $z\,\mathrm{pl.}$ est summa rectangulj sub dictis lateribus, et quadratj alterius lateris, quod æquatur potentiâ $G\,\mathrm{pl.}$ Et cæt.

\end{document}
